\chapter{Suplemento experimental}
\label{sec:supplement}

Para cada capítulo, reunimos aqui suas principais afirmações e aquilo em que elas se apoiam: o experimento em que foram medidas, o que exatamente foi medido e onde foi publicado. O status na coluna da direita diz o quanto a afirmação é firme: \emph{medido}, há um experimento direto; \emph{teoria}, uma consequência matemática deduzida no capítulo ou num trabalho clássico; \emph{modelo}, o resultado de um cálculo com um modelo do livro (scripts na pasta \texttt{models/}); \emph{estimativa}, um número de ordem de grandeza sem medição direta; \emph{hipótese}, uma explicação plausível que ainda não foi provada por experimento. Onde não há experimento, isso é dito.

\section{Capítulo~\ref{sec:neurontypes}. Tipos de neurônios}

Os números do capítulo se apoiam em contagens diretas de células no cérebro humano, onde elas existem; a regra ``um neurônio, um neurotransmissor'', em atlas celulares e em experimentos que também encontraram as exceções; o papel de \nt{DOPA}, em registros de disparos; e os papéis dos demais canais de difusão, em hipóteses.

{\footnotesize
\setlength{\tabcolsep}{3pt}\renewcommand{\arraystretch}{1.15}
\begin{longtable}{>{\raggedright}p{0.5cm}>{\raggedright}p{4.0cm}>{\raggedright}p{5.6cm}>{\raggedright}p{2.3cm}>{\raggedright\arraybackslash}p{1.7cm}}
\toprule
N.º & Proposição & Experimento e o que foi medido & Fonte & Status \\
\midrule\endfirsthead
\toprule
N.º & Proposição & Experimento e o que foi medido & Fonte & Status \\
\midrule\endhead
1 & O cérebro humano tem cerca de $8.6\cdot10^{10}$ neurônios, e quase todos os neurônios \nt{GLUT} são pequenos neurônios do cerebelo (tab.~\ref{tab:types}) & Fracionador isotrópico: o tecido é dissolvido numa suspensão homogênea de núcleos, e contam-se os núcleos com o marcador neuronal NeuN. Em homens adultos, $86.1\pm8.1$ bilhões de neurônios; apenas $19\,\%$ deles estão no córtex cerebral, a maior parte está no cerebelo & \cite{Azevedo2009} & medido \\
2 & Quase todo neurônio tem um neurotransmissor principal (proposição~\ref{prop:dale}), mas há exceções & Atlas transcriptômico do cérebro do camundongo: cerca de $4\cdot10^6$ células, $5322$ clusters; a cada cluster se atribui um neurotransmissor pelos genes de síntese e empacotamento. As exceções foram medidas diretamente: neurônios \nt{DOPA} em fatias liberam também GABA pelo mesmo transportador vesicular e inibem neurônios do estriado; em parte dos axônios \nt{DOPA}, glutamato e dopamina saem de terminais diferentes do mesmo fio (microscopia eletrônica e optogenética) & \cite{Strata1999,Yao2023,Tritsch2012,Zhang2015} & medido \\
3 & Toda a coluna da matriz de pesos tem o mesmo sinal~\eqref{eq:dalesign} & Dedução no capítulo a partir da proposição~\ref{prop:dale} e do fato de que os receptores de glutamato são quase todos excitatórios, e os de GABA, inibitórios. Não há verificação direta de ``todos os destinatários de um neurônio recebem peso do mesmo sinal'' como regra geral; contraexemplos: a liberação conjunta de GABA por neurônios \nt{DOPA} (linha~2) e destinatários com receptores de sinais diferentes (linha~11) & dedução no capítulo & teoria \\
4 & De três quartos a quatro quintos dos neurônios do córtex são \nt{GLUT}, de um quinto a um quarto são \nt{GABA} & Imunomarcação para GABA em colunas de $50$\,\textmu m de largura por toda a espessura do córtex em $10$ áreas corticais de macaco: em $8$ áreas os neurônios GABAérgicos são cerca de $25\,\%$ de todos os neurônios; no córtex visual primário, menos de $20\,\%$ & \cite{Hendry1987} & medido \\
5 & Um neurônio \nt{GLUT} tem cerca de $10^4$ saídas & Estereologia em autópsias: no neocórtex de homens jovens, $1.64\cdot10^{14}$ sinapses (microscopia eletrônica, dissector) e cerca de $2.3\cdot10^{10}$ neurônios. A razão dá $\sim7\cdot10^3$ sinapses por neurônio; em média, o número de entradas é igual ao de saídas & \cite{Tang2001synapses,Pakkenberg1997} & medido \\
6 & A saída de um neurônio \nt{DOPA} se ramifica em $10^5$--$10^6$ terminais; é uma estimativa pela cobertura do alvo, não uma contagem & Marcação viral da membrana de neurônios \nt{DOPA} isolados de rato e reconstrução completa do axônio: um neurônio cobre de $0.45$ a $5.7\,\%$ (em média $2.7\,\%$) do volume do estriado. Mediu-se a cobertura, não o número de terminais: este foi deduzido da cobertura em ordem de grandeza, sem contagem direta dos terminais. Para \nt{SERO} e \nt{NORA}, os números de terminais são estimativas grosseiras & \cite{Matsuda2009} & medido \\
7 & Os neurônios de difusão são poucos: \nt{DOPA} $\sim5\cdot10^5$ (o principal dos dois grupos, limite inferior), \nt{SERO} $\sim3\cdot10^5$ (soma de duas contagens parciais), \nt{HIST} $\sim6\cdot10^4$ (tab.~\ref{tab:types}, fig.~\ref{fig:typepie}) & Contagens estereológicas em humanos: $550\,000$ neurônios pigmentados na substância negra (sem a área tegmental ventral); $235\,000$ neurônios no núcleo dorsal da rafe (todos os neurônios do núcleo) e mais $125\,000$ neurônios serotoninérgicos no tegmento pontino; cerca de $64\,000$ neurônios histaminérgicos do hipotálamo. Para \nt{NORA}, \nt{GLYC}, \nt{ACET} e \nt{ADRE} não há contagem direta: na tabela são estimativas grosseiras de ordem de grandeza & \cite{Pakkenberg1991,Baker1990,Baker1991,Airaksinen1991} & medido \\
8 & O glutamato na fenda sobe até cerca de um milimolar e cai em $\sim1\ms$ & Pelo deslocamento de um bloqueador competitivo de dissociação rápida dos receptores NMDA durante a transmissão sináptica (cultura de neurônios do hipocampo): pico de $1.1$\,mM, constante de decaimento de $1.2\ms$ & \cite{Clements1992} & medido \\
9 & No córtex em funcionamento, as correntes excitatória e inibitória estão quase equilibradas, e o neurônio fica perto do limiar & Teoria: uma rede com conexões fortes chega sozinha ao regime equilibrado. Experimento: no córtex pré-frontal do furão in vivo, durante os estados ``up'', o potencial de reversão da corrente sináptica se mantém em torno de $-37$\,mV por centenas de milissegundos, embora a condutância varie em $\sim21$\,nS: excitação e inibição sobem e descem juntas & \cite{vanVreeswijk1996,Haider2006} & medido \\
10 & Os neurônios \nt{DOPA} informam o erro de predição de recompensa~\eqref{eq:rpe} (fig.~\ref{fig:rpe}) & Disparos de neurônios \nt{DOPA} de macaco: rajada ao suco inesperado; após o aprendizado, rajada ao sinal e nenhuma resposta ao suco previsto; queda da frequência quando o suco prometido não é dado. Num experimento quantitativo, a frequência é proporcional à diferença entre a recompensa atual e a média ponderada das anteriores, mas só para erros positivos. A equação~\eqref{eq:rpe} em si é o algoritmo de diferenças temporais & \cite{Schultz1997,Bayer2005,SuttonBarto2018} & medido \\
11 & O sinal e a velocidade são definidos pelo receptor: ionotrópicos, frações de milissegundo; metabotrópicos, muito mais lentos (proposição~\ref{prop:receptor}) & Pulsos rápidos de glutamato em fragmentos de membrana de neurônios do hipocampo: a corrente pelos receptores ionotrópicos sobe (de $20$ a $80\,\%$) em $0.2$--$0.6\ms$ e decai em $2$--$3\ms$. O potencial inibitório lento das células piramidais é abolido por um bloqueador seletivo do receptor metabotrópico de GABA. Duas famílias de receptores de dopamina com ação oposta; no estriado, os neurônios com receptor D1 precisam de dopamina para reforçar sinapses, e os com D2, para enfraquecê-las & \cite{Colquhoun1992,DutarNicoll1988,Beaulieu2011,Shen2008} & medido \\
12 & O canal de difusão não é um fio só, mas um feixe de poucas linhas (seção~\ref{sec:buses}) & Marcação viral e genética dos neurônios serotoninérgicos do núcleo dorsal da rafe do camundongo: os neurônios que enviam saída para a amígdala e para o córtex frontal ficam em lugares diferentes do núcleo, ramificam-se em regiões complementares e respondem de modo oposto a um estímulo aversivo. Revisão: subgrupos de neurônios do locus coeruleus (\nt{NORA}) codificam processos diferentes & \cite{Ren2018,Poe2020} & medido \\
13 & \nt{NORA} define o ganho $g$ & Hipótese do ganho adaptativo; modelo de rede em que as catecolaminas aumentam a inclinação da característica de transferência. Não há medição direta de ``$g$ cresce com a noradrenalina'' na rede inteira; mediu-se que o diâmetro da pupila acompanha os disparos dos neurônios do locus coeruleus do macaco & \cite{AstonJones2005,ServanSchreiber1990,Joshi2016} & hipótese \\
14 & \nt{ACET} alterna ``escrita/leitura'' e define a taxa de aprendizado & Hipótese. Base: em fatias de córtex olfatório de rato, agonistas da acetilcolina ($100$\,\textmu M) enfraquecem as sinapses internas, de ``lembrança'', em mais de $60\,\%$, e as de entrada em menos de $15\,\%$ & \cite{Hasselmo2006,HasselmoBower1992} & hipótese \\
15 & \nt{SERO} define o fator de desconto $\gamma$; os neurônios serotoninérgicos trabalham de forma regular, alguns disparos por segundo, e quase silenciam numa das fases do sono & Hipótese ``serotonina $=\gamma$''. O argumento mais forte: a ativação optogenética dos neurônios serotoninérgicos do núcleo dorsal da rafe do camundongo reduz o número de desistências prematuras na espera de uma recompensa adiada. A frequência de disparos e o silêncio no sono com movimentos oculares rápidos foram medidos (revisão de registros) & \cite{Doya2002,Miyazaki2014,JacobsAzmitia1992} & hipótese \\
\bottomrule
\end{longtable}}

\section{Capítulo~\ref{sec:glutcomputer}. O que os neurônios \nt{GLUT} calculam}

As afirmações lógicas do capítulo são teoremas sobre circuitos de pesos não negativos, deduzidos no próprio capítulo; a experiência confirma o modelo de neurônio e o fato de que circuitos montados segundo esses teoremas (detector de coincidência, linha de atraso, grupo autossustentado, cadeia) de fato existem no cérebro.

{\footnotesize
\setlength{\tabcolsep}{3pt}\renewcommand{\arraystretch}{1.15}
\begin{longtable}{>{\raggedright}p{0.5cm}>{\raggedright}p{4.0cm}>{\raggedright}p{5.6cm}>{\raggedright}p{2.3cm}>{\raggedright\arraybackslash}p{1.7cm}}
\toprule
N.º & Proposição & Experimento e o que foi medido & Fonte & Status \\
\midrule\endfirsthead
\toprule
N.º & Proposição & Experimento e o que foi medido & Fonte & Status \\
\midrule\endhead
1 & O neurônio é um circuito RC com limiar e reset~\eqref{eq:glutneuron} & No corpo de neurônios piramidais do córtex de rato (fatias) injetou-se uma corrente ruidosa, semelhante ao fluxo sináptico no cérebro vivo. A dependência da frequência média em relação à média e à dispersão da corrente é descrita pelo neurônio integrador com vazamento, se for acrescentada uma adaptação lenta da frequência. As faixas de $\tau$ e dos atrasos são dadas no capítulo sem referência experimental & \cite{Rauch2003} & medido \\
2 & O modelo~\eqref{eq:glutneuron} é uma simplificação: os ramos de entrada emitem eles mesmos disparos locais & Registro direto dos dendritos de neurônios piramidais do córtex visual do camundongo in vivo: o estímulo visual evoca disparos dendríticos locais que dependem dos receptores NMDA; suprimi-los alarga a sintonia do neurônio à orientação & \cite{Smith2013} & medido \\
3 & Janela de coincidência $\Delta_{\max}=\tau\ln\frac{w}{\theta-w}$~\eqref{eq:window}; com $\theta=1.5w$ ela vale $0.7\tau$ & Consequência do modelo~\eqref{eq:glutneuron}. Há medição direta de uma janela estreita de soma em neurônios com inibição rápida (capítulo~\ref{sec:gaba}, linha~3 de sua tabela) & dedução no capítulo & teoria \\
4 & Uma rede só de neurônios \nt{GLUT} calcula apenas funções monótonas das entradas (proposição~\ref{prop:monotone}) & Demonstração no capítulo: iteração a partir do silêncio com lado direito não decrescente. É uma afirmação sobre o modelo; não há experimento que a teste numa rede viva sem inibição & dedução no capítulo & teoria \\
5 & Os órgãos dos sentidos fornecem um par de fios: umas células respondem ao aumento do estímulo, outras à diminuição & Retina: já nas células bipolares, o sinal dos cones se divide em canais liga e desliga. Bloqueio farmacológico das bipolares liga no macaco: os canais seguem separados até o córtex; após o bloqueio piora a detecção do aumento de luz, mas não da diminuição & \cite{Schiller1992} & medido \\
6 & Com código de dois fios, uma rede de neurônios \nt{GLUT} calcula qualquer função lógica de modo assíncrono, sem ciclo comum, ao preço do dobro de neurônios (proposição~\ref{prop:dualrail}, \eqref{eq:dualrail}, fig.~\ref{fig:dualxor}) & O código de dois fios e a detecção de prontidão pelo próprio sinal são técnica padrão de circuitos assíncronos insensíveis a atrasos. O circuito de OU exclusivo com seis neurônios é montado no capítulo. Não há experimento mostrando que o cérebro calcula funções lógicas justamente com código de dois fios & \cite{Sparso2001} & teoria \\
7 & A maior diferença de tempo de chegada do som aos ouvidos no ser humano é $\approx0.6\ms$, e o cérebro distingue cerca de dez microssegundos & Ouvintes num experimento com duas alternativas de resposta: o limiar de discriminação da diferença de tempo (75\,\% de acertos) é de $9$\,\textmu s para ruído na banda de $150$--$1700$\,Hz, $11$\,\textmu s para um tom de $1$\,kHz, $28$\,\textmu s para um clique. O limite de $0.6\ms$ é cálculo do capítulo, $0.22\,\text{m}/343\,\text{m/s}$ & \cite{KlumppEady1956} & medido \\
8 & Nas aves, a diferença de tempos vira lugar por meio de linhas de atraso e detectores de coincidência~\eqref{eq:itd} (fig.~\ref{fig:itd}) & Suindara: os axônios dos dois ouvidos entram no núcleo dos detectores de coincidência por lados opostos; o tempo de chegada dos disparos ao longo do axônio muda de forma ordenada com a profundidade, e a variação dos atrasos através do núcleo ($160$\,\textmu s em $700$\,\textmu m) coincide com a faixa de atrasos interaurais & \cite{Carr1990} & medido \\
9 & Nos mamíferos, a fileira de atrasos não foi encontrada; a mesma tarefa é resolvida por detectores de coincidência com inibição de tempo preciso vinda de neurônios \nt{GLYC} & Registro in vivo na oliva superior medial do gerbil: as respostas não formam um mapa de direções; a aplicação de glicina e de seu bloqueador por micropipeta mostra que uma inibição \emph{glicinérgica} com tempo calculado com precisão define a faixa de trabalho dos atrasos. A inibição aqui vem de \nt{GLYC}, não de \nt{GABA} & \cite{Brand2002,Grothe2010} & medido \\
10 & Um grupo com forte realimentação é um flip-flop; a atividade autossustentada dura segundos após o estímulo (fig.~\ref{fig:bistable}) & Córtex pré-frontal do macaco numa tarefa de movimento ocular adiado para um ponto memorizado (atraso de $1$--$6$\,s): $87$ de $170$ neurônios ligados à tarefa mudam de frequência durante o atraso, e em $79\,\%$ deles isso depende da direção do ponto memorizado. Que essa atividade é mantida por uma realimentação com dois estados estáveis é teoria & \cite{Funahashi1989,Wang2001} & medido \\
11 & O bit é gravado se o aporte durar mais de $\approx0.7\tau$ (fig.~\ref{fig:recurrentgroup}b) & Cálculo da equação $\tau\,\dd r/\dd t=-r+f(wr+I)$ pelo método de Euler com $w=1$, $I=0.6$; não há script em \texttt{models/}. Não há experimento & cálculo no capítulo & modelo \\
12 & A memória pode ser guardada na facilitação das sinapses, quase sem disparos & Teoria: o cálcio residual nos terminais mantém o registro por cerca de um segundo. Base: no córtex pré-frontal medial do furão (registro de vários neurônios ao mesmo tempo) há uma sub-rede de neurônios piramidais ligados por sinapses facilitadoras com efeito prolongado. Revisão das observações de intervalos ``silenciosos'' durante a retenção na memória. Qual maneira o córtex usa e quando é questão em aberto & \cite{Mongillo2008,WangMarkram2006,Stokes2015} & hipótese \\
13 & A memória é apagada pela fadiga: o estoque de neurotransmissor se esgota & Registros pareados de neurônios piramidais do córtex: a resposta da sinapse cai com disparos frequentes, e a velocidade da queda é dada pela probabilidade de liberação. Que é exatamente isso que apaga o grupo autossustentado não foi mostrado diretamente & \cite{Tsodyks1997} & medido \\
14 & Uma cadeia de grupos é um temporizador disparado por evento; nas aves canoras, os neurônios disparam estritamente em sequência & Registro de neurônios identificados do centro vocal superior do mandarim durante o sono e o canto: cada neurônio que envia saída ao núcleo motor emite uma única rajada curta num momento preciso do canto, e os neurônios disparam em sequência & \cite{Hahnloser2002} & medido \\
\bottomrule
\end{longtable}}

\section{Capítulo~\ref{sec:gaba}. \nt{GLUT} e \nt{GABA} juntos}

As afirmações lógicas e dinâmicas do capítulo são deduzidas nele mesmo, por análise linear e pela lei de Ohm; a experiência mostra que os circuitos em que elas se apoiam (veto com atraso, inibição direta logo após a excitação, estabilização pela inibição, ritmo do laço \nt{GLUT}--\nt{GABA}) de fato funcionam no cérebro.

{\footnotesize
\setlength{\tabcolsep}{3pt}\renewcommand{\arraystretch}{1.15}
\begin{longtable}{>{\raggedright}p{0.5cm}>{\raggedright}p{4.0cm}>{\raggedright}p{5.6cm}>{\raggedright}p{2.3cm}>{\raggedright\arraybackslash}p{1.7cm}}
\toprule
N.º & Proposição & Experimento e o que foi medido & Fonte & Status \\
\midrule\endfirsthead
\toprule
N.º & Proposição & Experimento e o que foi medido & Fonte & Status \\
\midrule\endhead
1 & Os neurônios \nt{GABA} são de um quinto a um quarto dos neurônios do córtex & Imunomarcação para GABA em $10$ áreas do córtex de macaco: cerca de $25\,\%$ dos neurônios em $8$ áreas, menos de $20\,\%$ no córtex visual primário. Revisão da diversidade dos neurônios inibitórios do córtex & \cite{Hendry1987,Markram2004} & medido \\
2 & O que a inibição faz é decidido em grande parte pelo local de contato: dendrito, corpo, local de nascimento do disparo, terminal do fio de outro neurônio & Revisão: as classes de neurônios \nt{GABA} do córtex diferem pelo alvo no destinatário. Registro simultâneo do corpo e do dendrito de um neurônio piramidal: a inibição direta é muito mais forte no corpo do que nos dendritos. A inibição no terminal do fio de outro neurônio, na medula espinal, reduz a liberação de neurotransmissor de uma única linha de entrada (revisão das medições) & \cite{Markram2004,Pouille2001,Rudomin1999} & medido \\
3 & A troca de sinal por um intermediário custa um atraso, cerca de um milissegundo; a inibição que chega logo após a excitação estreita a janela de coincidência & Neurônios piramidais do hipocampo de rato em fatias: a inibição por um intermediário chega logo após a excitação direta com atraso curto, e a janela em que duas entradas excitatórias se somam até o limiar é menor que $2\ms$. Nos dendritos, onde essa inibição é mais fraca, a janela é mais larga & \cite{Pouille2001} & medido \\
4 & O veto $a\wedge\bar b$~\eqref{eq:veto} com OU e a constante um é funcionalmente completo (proposição~\ref{prop:gabacomplete}) & Demonstração no capítulo. Resultado clássico: uma rede de elementos de limiar com entradas inibitórias realiza qualquer expressão lógica. Afirmação sobre o modelo, sem experimento & \cite{McCullochPitts1943} & teoria \\
5 & O veto com atraso dá um detector de direção do movimento com velocidade ótima $v^*=\Delta x/d$~\eqref{eq:vstar} & Retina do coelho: a seletividade à direção é criada por uma inibição que, no movimento na direção ``nula'', anula a excitação. O registro separado das entradas excitatória e inibitória das células seletivas mostrou que as células amácrinas estreladas (\nt{GABA}) as inibem de forma assimétrica, só pelos prolongamentos voltados para o lado ``nulo''. A fórmula de $v^*$ é dedução do capítulo & \cite{Barlow1965,Fried2002} & medido \\
6 & A inibição por derivação divide a tensão em vez de subtrair~\eqref{eq:shunt} (fig.~\ref{fig:shunt}) & Lei de Ohm para um divisor de três condutâncias, com o potencial de equilíbrio do cloreto perto do repouso; para um neurônio sem disparos & dedução no capítulo & teoria \\
7 & Sobre a frequência de disparos a derivação age por subtração, e a divisão vem da inibição junto com ruído de fundo equilibrado & Modelo: num neurônio que dispara, a corrente pela derivação é quase constante, e o efeito sobre a frequência é subtrativo. Experimento: em neurônios piramidais do córtex somatossensorial de rato (fatias) introduziram-se, por fixação dinâmica, condutâncias de fundo excitatórias e inibitórias; o aumento simultâneo e equilibrado de ambas divide a inclinação da resposta sem deslocamento notável da frequência. Revisão: a divisão pela atividade total ocorre em muitas regiões do cérebro & \cite{HoltKoch1997,Chance2002,Carandini2012} & medido \\
8 & Com $\beta>1$ a inibição mútua escolhe um único vencedor (proposição~\ref{prop:wta}, \eqref{eq:wta}) & Os autovalores $-(1\pm\beta)/\tau$ são dedução do capítulo. As frequências $0.73$ e $0.53$ com $\beta=0.5$ são o ponto fixo $r_{1,2}=(I_{1,2}-\beta I_{2,1})/(1-\beta^2)$; não há script em \texttt{models/}. O capítulo não cita experimento direto & dedução no capítulo & teoria \\
9 & Reset da memória por sinal via \nt{GABA}; dois grupos com inibição mútua formam um latch & Decorre da fig.~\ref{fig:bistable} e da proposição~\ref{prop:wta}. Sem experimento & dedução no capítulo & teoria \\
10 & O equilíbrio do laço \nt{GLUT}--\nt{GABA} é estável se a inibição for forte e rápida o bastante (proposição~\ref{prop:ei}) & Modelo de duas populações~\eqref{eq:ei}; as condições (i) e (ii) são o determinante e o traço de sua matriz, deduzidas no capítulo & \cite{WilsonCowan1972} & teoria \\
11 & Com $w_{EE}=1.5$ e $w_{EI}w_{IE}=2$, a inibição rápida ($\tau_I=2\ms$) mantém $E^*=0.67h$, e a lenta ($30\ms$) provoca oscilações (fig.~\ref{fig:ei}) & Solução numérica de~\eqref{eq:ei}, script \texttt{figures/make\_ei.py} & cálculo & modelo \\
12 & O córtex trabalha no regime de estabilização pela inibição; a assinatura: empurre os neurônios \nt{GABA}, e sua frequência cai~\eqref{eq:isn} & A teoria previu a resposta paradoxal. Experimento: registros intracelulares no córtex visual primário do gato; quando a resposta é suprimida por um estímulo ao redor, a condutância inibitória não cresce, mas cai junto com a excitatória. A excitação optogenética de neurônios PV do camundongo nos córtices visual, somatossensorial e motor reduz a frequência dos neurônios inibitórios, inclusive sem estímulo e sob anestesia leve. O coeficiente $-1/3$ com os números da fig.~\ref{fig:ei} é dedução do capítulo & \cite{Tsodyks1997isn,Ozeki2009,Sanzeni2020} & medido \\
13 & O laço ``\nt{GLUT} excita \nt{GABA}, \nt{GABA} inibe \nt{GLUT}'' gera um ritmo rápido com período de $12$--$50\ms$ & A estimulação optogenética de neurônios \nt{GABA} rápidos no córtex somatossensorial do camundongo a $8$--$200$\,Hz amplifica seletivamente o ritmo gama ($20$--$80$\,Hz, período de $12$--$50\ms$); a estimulação de neurônios \nt{GLUT} amplifica só ritmos mais lentos & \cite{Cardin2009,Buzsaki2012} & medido \\
\bottomrule
\end{longtable}}

\section{Capítulo~\ref{sec:code}. Código Morse}

As estimativas-limite do capítulo são teoria da informação e contas; foram medidos onde surge o ruído no canal, quão rápido o cérebro trabalha e quanto custa um disparo, e a questão de que código o córtex usa de fato continua em aberto.

{\footnotesize
\setlength{\tabcolsep}{3pt}\renewcommand{\arraystretch}{1.15}
\begin{longtable}{>{\raggedright}p{0.5cm}>{\raggedright}p{4.0cm}>{\raggedright}p{5.6cm}>{\raggedright}p{2.3cm}>{\raggedright\arraybackslash}p{1.7cm}}
\toprule
N.º & Proposição & Experimento e o que foi medido & Fonte & Status \\
\midrule\endfirsthead
\toprule
N.º & Proposição & Experimento e o que foi medido & Fonte & Status \\
\midrule\endhead
1 & As frequências dos neurônios \nt{GLUT} do córtex vão de frações a dezenas de disparos por segundo, e a maior parte do tempo os neurônios trabalham perto do limite inferior & Registro com pipeta encostada à célula (a escolha do neurônio não depende de sua atividade) no córtex auditivo do rato acordado: a cada momento, menos de $5\,\%$ dos neurônios têm frequência acima de $20$ disparos por segundo; parte dos neurônios tem frequência de fundo abaixo de $0.01$ por segundo; a distribuição das frequências é lognormal & \cite{Hromadka2008} & medido \\
2 & A capacidade do canal de disparos é $R_{\max}\approx r\log_2\frac{e}{r\Delta t}$~\eqref{eq:spikeentropy}, e quase toda ela está no tempo dos disparos (proposição~\ref{prop:capacity}) & Entropia de uma cadeia binária de células de comprimento $\Delta t$. Números do capítulo: $8$ bits por disparo com $r=10$ por segundo e $\Delta t=1\ms$, cerca de $5$ bits com $\Delta t=10\ms$, menos de $2$ bits por segundo para um contador de disparos & \cite{MacKay1952} & teoria \\
3 & Os neurônios sensoriais transmitem de um a alguns bits por disparo, nos melhores casos até metade do limite & Reconstrução do estímulo a partir dos disparos de neurônios sensoriais para os quais o estímulo é conhecido; resumo das medições no livro & \cite{Rieke1997} & medido \\
4 & O gerador de disparos é preciso; o tempo exato é definido pelas variações rápidas da entrada & Neurônios do córtex de rato em fatias: com uma corrente repetida com flutuações rápidas, os disparos se repetem com precisão melhor que $1\ms$; com corrente constante, os instantes se espalham & \cite{Mainen1995} & medido \\
5 & A sinapse não é confiável: mais da metade dos disparos não chega ao destinatário por causa do acaso da liberação & Estimulação mínima das entradas de neurônios piramidais do hipocampo (fatias): mais da metade dos disparos não provoca resposta. Foram excluídas flutuações do limiar de estimulação, falhas de condução nas ramificações do axônio e a baixa temperatura do experimento; resta a liberação probabilística & \cite{AllenStevens1994} & medido \\
6 & O fluxo de disparos no córtex vivo parece aleatório: os intervalos se espalham como num fluxo de Poisson, e a variância do número de disparos fica próxima da média; parte do ``ruído'' é mensagem sobre os movimentos do animal & Registros nas áreas visuais V1 e MT do macaco acordado: coeficiente de variação dos intervalos de cerca de $1$, razão entre variância e média do número de disparos como num processo aleatório; um modelo que soma pequenas entradas independentes dá uma dispersão muito menor. No córtex visual do camundongo, uma parte considerável da variação é prevista pelo vídeo dos movimentos do próprio animal & \cite{Softky1993,Stringer2019} & medido \\
7 & O código de taxa é lento: erro $1/\sqrt{rT}$~\eqref{eq:rateerror}, enquanto o cérebro reconhece uma imagem em $\sim150\ms$ & A fórmula é estatística de Poisson, dedução do capítulo. Experimento: pessoas decidiam se havia um animal numa fotografia desconhecida mostrada por $20\ms$; os potenciais cerebrais para ``sim'' e ``não'' divergem após cerca de $150\ms$. A divisão desse tempo em uma dezena de estágios de $10$--$20\ms$ é estimativa do capítulo, não medição & \cite{Thorpe1996} & medido \\
8 & A média sobre um grupo é limitada pela correlação: o erro cai no máximo $1/\sqrt c$ vezes~\eqref{eq:correlated} (proposição~\ref{prop:pooling}) & Pares de neurônios da área MT do macaco registrados ao mesmo tempo: coeficiente de correlação dos números de disparos de $0.12$ em média, e a análise teórica mostra que mesmo essa correlação limita o ganho do grupo. Teoria: o limite é imposto só pela parte da correlação que se parece com o sinal. Modelo da posição oposta: a frequência de um grupo de $50$--$100$ neurônios é lida num único intervalo entre disparos, $10$--$50\ms$, e o tempo de disparos individuais não é usado no córtex & \cite{Zohary1994,MorenoBote2014,ShadlenNewsome1998} & medido \\
9 & O número pode ser escrito no tempo do primeiro disparo~\eqref{eq:latency}, na ordem dos disparos de um grupo ou na fase de um ritmo local & Retina da salamandra: a estrutura espacial de uma imagem mostrada brevemente está escrita nas latências relativas dos primeiros disparos dos neurônios de saída, quase independentemente do contraste. Células de lugar do hipocampo do rato: ao passar pelo ``seu'' lugar, os disparos se deslocam cada vez mais cedo no ciclo do ritmo teta ($7$--$12$\,Hz), um deslocamento de $100^\circ$ a $355^\circ$. Quanto o córtex usa o tempo dos disparos é questão em aberto (linha~8) & \cite{Gollisch2008,OKeefe1993} & medido \\
10 & Qual código é lido é decidido pela constante de tempo do destinatário~\eqref{eq:reader} (proposição~\ref{prop:reader}); no córtex ativo, os neurônios estão mais perto de detectores de coincidência & A fórmula~\eqref{eq:reader} é dedução do capítulo. Revisão de registros in vivo: no córtex ativo, a condutância da membrana é várias vezes maior que em repouso, e a constante de tempo, correspondentemente mais curta. Que o córtex troca de código exatamente assim não foi mostrado diretamente & \cite{Destexhe2003} & hipótese \\
11 & A sinapse com depressão transmite a mudança de frequência, essencialmente $\Delta r/r$~\eqref{eq:depression} (fig.~\ref{fig:depression}) & Registros pareados de neurônios piramidais do córtex: a velocidade da depressão é dada pela probabilidade de liberação; com depressão lenta, a resposta reflete a frequência; com rápida, as coincidências. Modelo baseado na depressão medida no córtex de rato: mudanças relativas iguais de frequência em entradas rápidas e lentas dão a mesma resposta, isto é, a sinapse transmite $\Delta r/r$. Os números $1.7\to2.2$, $6.7$ e $90\ms$ são cálculo do capítulo & \cite{Tsodyks1997,Abbott1997depression} & medido \\
12 & A rajada é o ``traço'': sua probabilidade de se perder, $(1-p)^k$~\eqref{eq:burst}, é pequena, e a facilitação a reduz ainda mais & Revisão: em sinapses pouco confiáveis, disparos isolados se perdem com frequência, e as rajadas são transmitidas de forma confiável graças à facilitação; a rajada provoca mudanças duradouras nas sinapses. Que o cérebro lê a rajada como um sinal à parte é hipótese & \cite{Lisman1997,AllenStevens1994} & hipótese \\
13 & Os neurônios \nt{GABA} rápidos definem o ritmo e a ``pontuação''; outra classe inibe os inibidores e abre a porta (proposição~\ref{prop:punctuation}) & Revisão das propriedades das classes de neurônios \nt{GABA} do córtex. Optogenética: a excitação rítmica de neurônios \nt{GABA} rápidos provoca o ritmo gama, e a resposta a um estímulo depende da fase do ciclo. No córtex visual do camundongo, os neurônios PV inibem uns aos outros, os SST inibem todas as outras classes, e os VIP inibem sobretudo os SST. A excitação dos neurônios VIP no camundongo acordado remove a inibição dos neurônios \nt{GLUT}; eles são ligados por um sinal de reforço. A divisão de papéis como regra geral é hipótese & \cite{Tremblay2016,Cardin2009,Pfeffer2013,Pi2013} & hipótese \\
14 & A energia limita a frequência média a algo da ordem de um disparo por segundo~\eqref{eq:energy} & Orçamento energético da substância cinzenta de roedores a partir de dados anatômicos e fisiológicos: disparos e correntes sinápticas respondem por $47\,\%$ e $34\,\%$ da energia de sinalização; no máximo $\sim15\,\%$ dos neurônios podem estar ativos ao mesmo tempo. Para o córtex humano: no máximo $\sim1\,\%$ dos neurônios podem estar fortemente ativos ao mesmo tempo. O número de $\sim10^9$ ATP por disparo e a potência de $\sim10$\,W para sinalização são estimativa do capítulo com base nesses cálculos & \cite{Attwell2001,Lennie2003} & teoria \\
15 & O código é esparso e ajustado para o máximo de bits por joule; o córtex troca precisão por energia (proposição~\ref{prop:economy}) & Hipótese do código econômico. Base: em camundongos com restrição alimentar, no córtex visual caem a condutância dos receptores AMPA e o gasto sináptico de ATP ($-29\,\%$), a frequência de disparos se mantém, e a sintonia à orientação se alarga $32\,\%$ & \cite{Barlow1961,Padamsey2022} & hipótese \\
\bottomrule
\end{longtable}}

\section{Capítulo~\ref{sec:serocomputer}. Neurônios \nt{SERO}}

As propriedades dos próprios neurônios \nt{SERO} (ritmo, autoinibição, sinal definido pelo receptor, subsistemas) foram medidas diretamente; os cálculos do capítulo (regulador, filtro, lógica) decorrem de suas equações; $\tau_c$, o coeficiente de difusão da serotonina e o número de locais de liberação são estimativas; o papel do laço como regulador de nível no cérebro é uma hipótese (no núcleo da rafe a autoinibição é pontual e dá só pausas curtas), e o papel do \nt{SERO} como horizonte é uma hipótese com experimentos comportamentais a seu favor.

{\footnotesize
\setlength{\tabcolsep}{3pt}\renewcommand{\arraystretch}{1.15}
\begin{longtable}{>{\raggedright}p{0.5cm}>{\raggedright}p{4.0cm}>{\raggedright}p{5.6cm}>{\raggedright}p{2.3cm}>{\raggedright\arraybackslash}p{1.7cm}}
\toprule
N.º & Proposição & Experimento e o que foi medido & Fonte & Status \\
\midrule\endfirsthead
\toprule
N.º & Proposição & Experimento e o que foi medido & Fonte & Status \\
\midrule\endhead
1 & O sinal da ação da serotonina é escolhido pelo receptor do destinatário (proposição~\ref{prop:receptor}) & Os receptores de serotonina se dividem em famílias por farmacologia e mecanismo: o 5-HT$_{1A}$, via proteína G, abre canais de potássio e inibe a célula; o 5-HT$_{2A}$ e o ionotrópico 5-HT$_3$ excitam. Um mesmo neurotransmissor produz respostas opostas em células diferentes. & \cite{BarnesSharp1999} & medido \\
2 & Na vigília, o neurônio \nt{SERO} emite regularmente alguns disparos por segundo & Registro de neurônios do núcleo dorsal da rafe em gatos em livre movimento: descarga rítmica lenta de $2.8\pm0.2$ disparos/s na vigília tranquila; no sono de ondas lentas a taxa cai $34$--$68\,\%$, no sono REM, $98\,\%$. Em ratos anestesiados, $1.7\pm0.5$ disparos/s, muito regular. & \cite{TrulsonJacobs1979, GagnonParent2014, JacobsAzmitia1992} & medido \\
3 & O neurônio \nt{SERO} é um marca-passo: não precisa de empurrão a cada disparo, mas precisa de um aporte de fundo constante (na fatia, noradrenalina via receptores $\alpha_1$; no organismo, do \nt{NORA}) & Na fatia de cérebro, as células descarregam de forma lenta e regular; após cada disparo há uma pós-hiperpolarização de $200$--$800\ms$. Na fatia há menos células com atividade espontânea que no organismo: o ritmo regular precisa de uma entrada tônica de noradrenalina via receptores $\alpha_1$, cujo bloqueio o apaga. & \cite{VandermaelenAghajanian1983} & medido \\
4 & Quase todos os receptores são metabotrópicos; a resposta é lenta: sobe e desce em cerca de um segundo & Todas as famílias de receptores, exceto a 5-HT$_3$, são acopladas a proteínas G. Na fatia, a liberação evocada de serotonina produz, via 5-HT$_{1A}$, uma corrente inibitória que sobe e desce em um segundo. & \cite{BarnesSharp1999, CourtneyFord2016} & medido \\
5 & Nas regiões de destino, a serotonina sai no volume, e não só na fenda; no próprio núcleo da rafe, a inibição dos vizinhos é ponto a ponto & Voltametria cíclica rápida em fatias: a liberação por disparo é igual numa região com sinapses e no núcleo da rafe, onde quase não há sinapses; não depende do bloqueio da recaptação e dos receptores; há muito menos moléculas por terminal que transportadores e receptores, e a concentração extracelular coincide com a afinidade do receptor principal. No núcleo da rafe, a inibição via 5-HT$_{1A}$ é ponto a ponto: o transportador não deixa a serotonina acumular fora da fenda. & \cite{Bunin1998, CourtneyFord2016} & medido \\
6 & A concentração em~\eqref{eq:seroneuron} diminui pela recaptação; $\tau_c$ da ordem de um segundo é uma estimativa & Voltametria no cérebro vivo: a serotonina extracelular é limitada sobretudo pela recaptação e pela degradação, não pela síntese. Não há medição direta de $\tau_c$ nos trabalhos citados; a ordem de grandeza é estimativa do capítulo. & \cite{Hashemi2012serotonin} & medido; $\tau_c$: estimativa \\
7 & NÃO e NOU num único marca-passo; completude da lógica \nt{SERO} (proposição~\ref{prop:serocomplete}, fig.~\ref{fig:serogates}) & Decorre de~\eqref{eq:seroneuron}: um marca-passo inibível é um NOU, e o NOU é funcionalmente completo. Não há experimento em que se tenha montado lógica com neurônios \nt{SERO}; a condição de volumes separados não se cumpre no cérebro. & dedução no capítulo & teoria \\
8 & A serotonina inibe o próprio neurônio \nt{SERO} por receptores nele & Em ratos anestesiados, agonistas de 5-HT$_{1A}$ por via intravenosa e iontoforética inibem a descarga dos neurônios da rafe dorsal com a mesma relação dose-resposta da própria serotonina; agonistas de 5-HT$_{1B}$ quase não agem. Na fatia, a serotonina endógena das células vizinhas produz, via 5-HT$_{1A}$, uma corrente e uma pausa curta na descarga, sem mudar a taxa média. & \cite{Sprouse1987, CourtneyFord2016} & medido \\
9 & No modelo, o laço por um volume comum é um regulador de nível: a dispersão e a constante de tempo caem $1+G$ vezes~\eqref{eq:feedback}; que o laço funcione assim no cérebro é hipótese & Fórmula clássica do amplificador com realimentação negativa; deduzida de novo no capítulo. O ganho de laço $G$ do sistema \nt{SERO} não foi medido. Medido: na fatia, a autoinibição é ponto a ponto, dá uma pausa curta e não muda a taxa média. & \cite{Black1934feedback, CourtneyFord2016} & teoria; no cérebro, hipótese \\
10 & A concentração é uma média móvel da taxa, um filtro passa-baixas~\eqref{eq:lowpass}, fig.~\ref{fig:lowpass} & Solução da equação linear~\eqref{eq:seroneuron}; a figura é um cálculo por essa fórmula com $\tau_c=1\,\text{s}$. Não há modelo separado em \texttt{models/}. & dedução no capítulo & teoria \\
11 & A tortuosidade dos espaços torna a difusão de uma molécula pequena cerca de $2.6$ vezes mais lenta & Método da fonte pontual: iontoforese de tetrametilamônio e registro de sua concentração com eletrodo íon-seletivo em fatias e no cérebro vivo. Tortuosidade $\lambda\approx1.6$, ou seja, o $D$ efetivo é $\lambda^2\approx2.6$ vezes menor que o livre; fração de volume extracelular de cerca de $0.2$. O coeficiente de difusão da própria serotonina no tecido não foi medido nesses trabalhos; no capítulo, é uma estimativa. & \cite{Sykova2008} & medido \\
12 & Comprimento de um ``fio'' independente $\ell\sim\sqrt{D\tau}\approx20$~\textmu m~\eqref{eq:diffusionlength}; o número de fios é limitado pelo número dessas regiões (proposição~\ref{prop:volumewires}) & Lei da difusão e substituição das estimativas de $D$ e $\tau$ (linhas~6 e~11); a proposição~\ref{prop:volumewires} é uma consequência. Não há experimento direto que conte os canais independentes da transmissão por volume. & dedução no capítulo & teoria \\
13 & A saída de um neurônio \nt{SERO} se espalha por enormes áreas do cérebro; locais de liberação, por estimativa, $\sim10^5$ & Reconstrução completa de neurônios isolados da rafe dorsal no rato: todos os axônios ascendentes se ramificam muito, o comprimento total do axônio de uma célula chega a $18.7$~cm, e os ramos de uma única célula alcançam o estriado, o córtex pré-frontal e a amígdala. O número de locais de liberação por célula não foi contado diretamente. & \cite{GagnonParent2014} & medido; $10^5$: estimativa \\
14 & Os neurônios \nt{SERO} se dividem em alguns subsistemas com saídas e respostas diferentes & Marcação viral-genética em camundongos: os neurônios que vão para a amígdala e para o córtex frontal quase não se sobrepõem na ramificação, recebem entradas diferentes e respondem de modo oposto a um estímulo aversivo; ativar e desativar cada grupo muda o comportamento de modos diferentes. & \cite{Ren2018} & medido \\
15 & Condição de paciência $D<\tau_\gamma\ln(R/r_0)$~\eqref{eq:patience} com desconto exponencial; com o desconto medido, hiperbólico, $D<\tau_\gamma(R/r_0-1)$ & Decorre do desconto $e^{-D/\tau_\gamma}$ ou $1/(1+D/\tau_\gamma)$; dedução no capítulo. Macacos, escolha com atrasos de segundos: o desconto fica mais perto da hipérbole que da exponencial; a resposta dos neurônios \nt{DOPA} ao sinal antecipador de uma recompensa adiada também cai com o atraso de forma hiperbólica. & dedução no capítulo; \cite{KobayashiSchultz2008} & teoria \\
16 & O \nt{SERO} define o horizonte $\tau_\gamma$ (seção~\ref{sec:horizon}) & A ativação optogenética dos neurônios \nt{SERO} da rafe dorsal em camundongos prolonga a espera por uma recompensa adiada e aumenta a persistência na busca de uma recompensa que ficou rara. Em humanos, a redução da serotonina (dieta sem triptofano) torna mais íngreme o desconto da recompensa adiada. Como a concentração se converte em $\tau_\gamma$ não se sabe. & \cite{Doya2002, Miyazaki2014, Lottem2018, Schweighofer2008} & hipótese \\
\bottomrule
\end{longtable}}

\section{Capítulo~\ref{sec:dofa}. Aprendizado com \nt{DOPA}}

Os mecanismos da sinapse (porções, detector de coincidência, cálcio, janelas de plasticidade) e o comportamento da dopamina como erro de predição foram medidos diretamente; que as regras levam ao gradiente e quanto custa a difusão são teoremas; o resultado do aprendizado de escolha é um cálculo com o modelo do livro; o circuito que calcula o erro é uma hipótese.

{\footnotesize
\setlength{\tabcolsep}{3pt}\renewcommand{\arraystretch}{1.15}
\begin{longtable}{>{\raggedright}p{0.5cm}>{\raggedright}p{4.0cm}>{\raggedright}p{5.6cm}>{\raggedright}p{2.3cm}>{\raggedright\arraybackslash}p{1.7cm}}
\toprule
N.º & Proposição & Experimento e o que foi medido & Fonte & Status \\
\midrule\endfirsthead
\toprule
N.º & Proposição & Experimento e o que foi medido & Fonte & Status \\
\midrule\endhead
1 & A retropropagação do erro, na forma que tem nas redes artificiais, não foi encontrada no cérebro & Não há experimento direto: é uma afirmação de ausência. As revisões analisam o que o algoritmo exige (conexões de retorno que repetem as diretas, fase separada) e que aproximações biológicas foram propostas. & \cite{Lillicrap2020, WhittingtonBogacz2019} & hipótese \\
2 & O peso se decompõe em número de locais de liberação, probabilidade de liberação e resposta a uma porção: $w=N_s\,p\,A_1$~\eqref{eq:quantal} & Junção neuromuscular de rã com liberação reduzida: a resposta do destinatário varia em degraus múltiplos da resposta miniatura espontânea a uma porção; o número de porções por disparo segue a lei de Poisson. & \cite{delCastillo1954} & medido \\
3 & O receptor do destinatário é um detector de coincidência; o cálcio desencadeia a mudança de peso & Patch-clamp de neurônios de camundongo: íons Mg$^{2+}$ bloqueiam o canal aberto pelo glutamato no potencial de repouso e saem com a despolarização. Fatias de hipocampo: um quelante de cálcio no destinatário bloqueia a potenciação de longo prazo, e a liberação de cálcio por luz dentro da célula por si só reforça a transmissão. & \cite{Nowak1984, Malenka1988calcium} & medido \\
4 & Qualquer lado executa a decisão: cresce $A_1$ no destinatário, muda $p$ no remetente & Glutamato liberado por luz sobre uma única espinha causa um crescimento rápido e duradouro da espinha e da corrente por seus receptores; a potenciação é acompanhada da inserção de receptores AMPA. A despolarização do destinatário suprime temporariamente a liberação no remetente; o efeito é levado por endocanabinoides que voltam pela fenda (mostrado em entradas \nt{GABA}). A depressão de curto prazo depende da probabilidade de liberação no remetente. & \cite{Matsuzaki2004, Malinow2002, WilsonNicoll2001, Tsodyks1997} & medido \\
5 & A sinapse se reforça quando remetente e destinatário estão ativos juntos~\eqref{eq:hebb} & Coelho anestesiado: após séries curtas de disparos na via de entrada do hipocampo, a resposta fica reforçada por $30$~min a $10$~h. Em fatia de hipocampo, os disparos do destinatário reforçam só as sinapses ativas no mesmo momento. & \cite{Bliss1973LTP, Kelso1986hebb} & medido \\
6 & A regra~\eqref{eq:oja} encontra o autovetor principal das correlações das entradas (proposição~\ref{prop:oja}) & Teorema; demonstração no capítulo. Não há experimento em que um neurônio tenha aprendido a componente principal de suas entradas; o mecanismo que limita o crescimento dos pesos na sinapse não foi estabelecido. & \cite{Oja1982} & teoria \\
7 & Hebb temporal: janela de cerca de $\pm20\ms$ (fig.~\ref{fig:stdp}); sequências repetidas fazem crescer cadeias & Pares de neurônios de hipocampo em cultura: um disparo do destinatário até $20\ms$ depois do disparo do remetente produz reforço duradouro, até $20\ms$ antes, enfraquecimento; ambos dependem de receptores NMDA. A razão $A_-=0.6A_+$ na figura é ilustrativa. Que assim crescem cadeias na rede não foi medido diretamente. & \cite{BiPoo1998} & medido \\
8 & Traço~\eqref{eq:threefactor}: a dopamina que chega $1$--$2\,\text{s}$ após a atividade conjunta reforça a conexão; mais tarde, não (proposição~\ref{prop:threefactor}) & Fatias de estriado, estimulação optogenética separada das entradas glutamatérgicas e dopaminérgicas: a espinha só cresce se a dopamina chega $0.3$--$2\,\text{s}$ após a entrada glutamatérgica; a janela é dada pela subida e queda rápidas do AMPc. No córtex, o pareamento deixa traços separados para reforço e enfraquecimento, que receptores de monoaminas convertem em mudanças duradouras numa janela da ordem de segundos. & \cite{Yagishita2014, He2015eligibility} & medido \\
9 & Janelas de segundos também existem sem dopamina: o terceiro sinal é uma forte excitação do destinatário & Camundongo vivo, campo CA1: um único evento de platô no destinatário reforça as entradas chegadas segundos antes e depois dele e cria um campo de lugar numa só passagem. & \cite{Bittner2017} & medido \\
10 & O passo médio da regra de três fatores segue o gradiente da recompensa~\eqref{eq:perturbation} (proposição~\ref{prop:gradient}) & Teorema do gradiente para aprendizado por desvios aleatórios; deduzido no capítulo para ruído pequeno. & \cite{Williams1992} & teoria \\
11 & O ruído é busca: a rede aprende com seus próprios desvios aleatórios & Mandarins jovens: desligar o núcleo de saída do circuito dos gânglios da base elimina, de forma abrupta e reversível, a variabilidade do canto. Mandarins adultos: uma interferência é aplicada às execuções de uma sílaba com altura de um lado da média, e as aves deslocam rapidamente a altura para o outro lado, aprendendo com a própria dispersão. & \cite{Olveczky2005, TumerBrainard2007} & medido \\
12 & A escolha entre duas opções é aprendida com $\tau_e=1\,\text{s}$ e $\delta=R-\bar R$; não com $\tau_e=50\ms$; sem a subtração, os pesos crescem sem limite, e com taxa de aprendizado alta a rede pode fixar a pior escolha & \texttt{models/dofa\_learning.py}, $\eta=0.05$, $500$ tentativas, $6$ sementes. $\tau_e=1\,\text{s}$: fração de escolhas de $A$ $0.65$--$0.79\to0.96$--$1.00$, pesos $0.85$--$1.33$. $\tau_e=50\ms$: $0.38$--$0.55$, os pesos não mudam. $\delta=R$: peso do vencedor $860$--$1190$. $\delta=R$, $\eta=0.5$, $3$ sementes: uma se fixou em $B$ ($0.04\to0$). O script imprime os quatro casos; o recálculo coincidiu com o capítulo. & \texttt{models/\allowbreak dofa\_\allowbreak learning.py} & modelo \\
13 & O tempo de aprendizado por um único sinal comum cresce proporcionalmente ao número de neurônios ruidosos (proposição~\ref{prop:broadcastcost}) & Curvas de aprendizado de uma rede linear: a perturbação de neurônios é mais lenta que o gradiente exato tantas vezes quantos neurônios de saída houver; a perturbação de pesos, ainda tantas vezes quantas entradas houver. & \cite{Werfel2005} & teoria \\
14 & As saídas do \nt{DOPA} vão sobretudo às regiões de escolha de ações & Reconstrução completa dos axônios de neurônios \nt{DOPA} nigroestriatais isolados do rato: os ramos de uma célula cobrem $0.45$--$5.7\,\%$ do volume do estriado (em média $2.7\,\%$); fora do estriado quase não há ramos. & \cite{Matsuda2009} & medido \\
15 & O córtex aprende a prever sua entrada, e o erro é calculado no local & Revisão: no córtex sensorial há respostas ao desacordo entre a entrada esperada e a recebida. Que isso seja uma regra geral de aprendizado do córtex não foi demonstrado. & \cite{Keller2018} & hipótese \\
16 & A dopamina se comporta como o erro~\eqref{eq:ctd}: pico, transferência para o sinal antecipador, queda (fig.~\ref{fig:ctdcases}) & Neurônios \nt{DOPA} de macaco: pico com suco inesperado; após o aprendizado, pico no sinal antecipador e nenhuma resposta ao suco prometido; queda quando o suco não veio. A forma contínua~\eqref{eq:ctd} é teoria. & \cite{Schultz1997, Doya2000} & medido \\
17 & Circuito que calcula $\dd V/\dd t$ e subtrai a expectativa & Camundongos, registro e optogenética na área tegmental ventral: os neurônios \nt{DOPA} subtraem a expectativa da recompensa; neurônios \nt{GABA} vizinhos os inibem quando a recompensa é esperada, e excitá-los ou inibi-los muda o erro. Como a derivada é calculada não foi mostrado. & \cite{Eshel2015arithmetic} & hipótese \\
18 & O neurônio \nt{DOPA} é um marca-passo; as quedas são menores que os picos & Em fatia, os neurônios \nt{DOPA} descarregam espontaneamente e de modo muito regular, sobre um potencial marca-passo lento. No macaco, a taxa acompanha quantitativamente o erro positivo e transmite o negativo apenas fracamente. & \cite{GraceOnn1989, Bayer2005} & medido \\
19 & Ator e crítico (fig.~\ref{fig:actorcritic}) & O algoritmo vem da teoria do aprendizado por reforço. Em humanos (fMRI), o estriado ventral se comporta como crítico com e sem escolha; o dorsal, só quando é preciso escolher ações. & \cite{SuttonBarto2018, ODoherty2004actor} & hipótese \\
20 & O sinal de $\delta$ na regra é invertido nos neurônios com a segunda família de receptores & Fatias de estriado: nos neurônios com receptores D1, o pareamento produz reforço, que requer D1; nos neurônios com D2, enfraquecimento, que requer D2. & \cite{Shen2008} & medido \\
\bottomrule
\end{longtable}}

\section{Capítulo~\ref{sec:dofaalt}. Outras teorias do \nt{DOPA}}

Cada um dos quadros ampliados se apoia em suas próprias medições diretas da dopamina (dispersão dos pontos de reversão, respostas ao desagradável e ao novo, crescimento lento, respostas à troca de sabor), mas as fórmulas que as explicam são teorias, e entre duas delas os experimentos por enquanto se contradizem.

{\footnotesize
\setlength{\tabcolsep}{3pt}\renewcommand{\arraystretch}{1.15}
\begin{longtable}{>{\raggedright}p{0.5cm}>{\raggedright}p{4.0cm}>{\raggedright}p{5.6cm}>{\raggedright}p{2.3cm}>{\raggedright\arraybackslash}p{1.7cm}}
\toprule
N.º & Proposição & Experimento e o que foi medido & Fonte & Status \\
\midrule\endfirsthead
\toprule
N.º & Proposição & Experimento e o que foi medido & Fonte & Status \\
\midrule\endhead
1 & A regra assimétrica~\eqref{eq:asymmetric} converge para o ponto~\eqref{eq:expectile}; com $\kappa=1/2$ é a média; para uma gota com probabilidade $1/2$, $V=\kappa$ (fig.~\ref{fig:expectiles}) & O ponto~\eqref{eq:expectile} é um expectil, solução de um problema de mínimos quadrados com pesos diferentes para desvios positivos e negativos; para o exemplo do capítulo, a dedução está no próprio capítulo. & \cite{NeweyPowell1987}; dedução no capítulo & teoria \\
2 & Neurônios \nt{DOPA} diferentes têm pontos de reversão diferentes, ordenados pela assimetria (proposição~\ref{prop:expectile}) & Camundongos, neurônios \nt{DOPA} da área tegmental ventral marcados optogeneticamente, recompensas de tamanhos diferentes: o ponto em que o pico dá lugar à queda difere entre neurônios; nos neurônios com respostas mais assimétricas ele é mais alto; das respostas do grupo se reconstrói a forma da distribuição das recompensas. Que essa seja a função principal da dispersão é hipótese. & \cite{Dabney2020} & medido \\
3 & Parte dos neurônios \nt{DOPA} responde ao desagradável com um pico & Macacos: alguns neurônios \nt{DOPA} são excitados pelo sinal de recompensa e inibidos pelo sinal de um jato de ar no olho; outros são excitados por ambos; os grupos ficam em partes diferentes do mesencéfalo. & \cite{Matsumoto2009, BrombergMartin2010} & medido \\
4 & O módulo do erro ou a novidade define a taxa de aprendizado & Nos trabalhos citados não há experimento direto em que o sinal de importância dos neurônios \nt{DOPA} mude a taxa de aprendizado; a revisão propõe o papel de sinal ``atenção, é importante''. & \cite{BrombergMartin2010} & hipótese \\
5 & Um fio separado para o eixo do perigo & Camundongos: os neurônios \nt{DOPA} que vão à extremidade posterior do estriado respondem a estímulos novos e fortes, e não à recompensa; sua destruição enfraquece a evitação de um objeto novo. & \cite{Menegas2018} & medido \\
6 & Tempo ótimo de ação $\tau_a^*=\sqrt{C/\bar R}$~\eqref{eq:vigor} & Mínimo da soma $C/\tau_a+\bar R\tau_a$; dedução no capítulo. No modelo original, a velocidade das ações é dada pelo preço do tempo, igual à recompensa média. & \cite{Niv2007}; dedução no capítulo & teoria \\
7 & O nível lento de dopamina leva $\bar R$ e define a velocidade das ações (proposição~\ref{prop:multiplex}) & Ratos, microdiálise e voltametria no núcleo accumbens: o nível de dopamina, minuto a minuto, acompanha a frequência das recompensas e a velocidade das ações. Em humanos, a L-DOPA reforça a dependência do tempo de reação em relação à frequência média de recompensas. Que o nível lento seja exatamente $\bar R$ não foi demonstrado. & \cite{Hamid2016, Beierholm2013vigor} & hipótese \\
8 & O nível lento vem de duas fontes: a liberação muda nas saídas sem mudança na taxa de disparos, mas o crescimento lento também aparece na própria taxa & Ratos, uma mesma tarefa: a liberação no núcleo accumbens acompanha a expectativa de recompensa em mudança, mas a taxa de disparos dos neurônios \nt{DOPA} identificados da área tegmental ventral não. Camundongos num corredor virtual (linha~10): o crescimento suave ao se aproximar da recompensa aparece também nos disparos dos neurônios \nt{DOPA} identificados. & \cite{Mohebi2019, Kim2020} & medido \\
9 & A dopamina sobe suavemente enquanto o animal se aproxima de uma recompensa distante & Ratos num labirinto, voltametria no estriado: a concentração de dopamina sobe em segundos à medida que o alvo se aproxima; a inclinação depende da distância e do tamanho da recompensa. & \cite{Howe2013ramp, Hamid2016} & medido \\
10 & O crescimento é o erro $\delta$ com refinamento da estimativa~\eqref{eq:ramp}; um salto no corredor dá um pico (fig.~\ref{fig:ramp}) & A fórmula e o número $\delta=0.75\,V$ por segundo são dedução do capítulo. Experimento: camundongos num corredor virtual, transporte instantâneo para mais perto do alvo; os disparos dos corpos celulares, o cálcio dos corpos e das saídas e a concentração no estriado respondem como erro, e não como a expectativa $V$. Qual das duas explicações vale em todas as saídas está em discussão. & \cite{Kim2020} & medido \\
11 & Olhar para trás: a dopamina informa uma medida de ligação causal~\eqref{eq:retro} & Camundongos, liberação de dopamina no núcleo accumbens em experimentos em que essa teoria e o erro~\eqref{eq:ctd} preveem coisas diferentes: em vários experimentos a liberação seguiu a primeira. & \cite{Jeong2022} & hipótese \\
12 & A resposta da dopamina durante o aprendizado se desloca gradualmente da recompensa para o sinal antecipador & Camundongos, sinal das saídas dos neurônios \nt{DOPA} no estriado ventral ao longo do aprendizado: o pico passa gradualmente do instante da recompensa para o do sinal antecipador, como exige o aprendizado por~\eqref{eq:ctd}. & \cite{Amo2022} & medido \\
13 & Os neurônios \nt{DOPA} respondem à troca do sabor da recompensa com o mesmo valor & Ratos, registro de neurônios da área tegmental ventral: trocar um suco por outro de valor igual provoca resposta, embora o erro~\eqref{eq:ctd} seja zero. & \cite{Takahashi2017} & medido \\
14 & Picos artificiais ensinam a ligação entre dois estímulos neutros & Ratos, experimento de pré-condicionamento sensorial de dois estímulos sem recompensa: a ativação optogenética dos neurônios \nt{DOPA} na transição do primeiro para o segundo cria a ligação, e a desativação impede que seja aprendida. & \cite{Sharpe2017} & medido \\
15 & Vetor de erros por característica~\eqref{eq:vectortd} & Teoria do erro de predição de características; não há verificação direta da forma vetorial. & \cite{Gardner2018} & hipótese \\
16 & Neurônios \nt{DOPA} diferentes numa mesma tarefa levam grandezas diferentes & Camundongos num corredor virtual, imagem de dois fótons de neurônios \nt{DOPA}: uns ligados à recompensa, outros à velocidade e à rotação, outros aos sinais antecipadores e à precisão da escolha. & \cite{Engelhard2019} & medido \\
17 & Um feixe em vez de um fio (proposição~\ref{prop:bundle}) & Resumo das linhas~2--16: cada decomposição foi medida separadamente; não há experimento que mostre um quadro único em que todas sejam partes de um mesmo sinal. & --- & hipótese \\
\bottomrule
\end{longtable}}

\section{Capítulo~\ref{sec:nora}. \nt{NORA}}

A organização do sistema \nt{NORA}, seus dois modos, a ligação com a pupila (não só via \nt{NORA}), o aumento da relação sinal-fundo e o U invertido no comportamento foram medidos diretamente; o ganho multiplicativo $g$ é um modelo; que o ganho comute o modo de escolha e funcione como inverso da temperatura são deduções do capítulo; a vantagem do ajuste do ganho é um cálculo com o modelo do livro; ``o \nt{NORA} é o botão da busca'' e ``o \nt{NORA} é uma interrupção'' são hipóteses.

{\footnotesize
\setlength{\tabcolsep}{3pt}\renewcommand{\arraystretch}{1.15}
\begin{longtable}{>{\raggedright}p{0.5cm}>{\raggedright}p{4.0cm}>{\raggedright}p{5.6cm}>{\raggedright}p{2.3cm}>{\raggedright\arraybackslash}p{1.7cm}}
\toprule
N.º & Proposição & Experimento e o que foi medido & Fonte & Status \\
\midrule\endfirsthead
\toprule
N.º & Proposição & Experimento e o que foi medido & Fonte & Status \\
\midrule\endhead
1 & Os neurônios \nt{NORA} no ser humano são algumas dezenas de milhares, quase todos num único grupo & Cortes seriados do tronco encefálico humano, coloração para tirosina hidroxilase: $53\,900$ células no locus coeruleus e mais $6\,260$ num subnúcleo vizinho. & \cite{Baker1989LC} & medido \\
2 & As saídas se espalham por quase todo o cérebro, mas partes do grupo atendem em parte regiões diferentes & Marcação viral-genética em camundongos: os neurônios \nt{NORA} do locus coeruleus recebem entradas de amplas regiões do cérebro e enviam saídas a amplas regiões do cérebro e da medula espinhal. Em ratos: o grupo que vai à amígdala reforça o aprendizado do medo, e um grupo separado que vai ao córtex pré-frontal, sua extinção. & \cite{Schwarz2015LC, Uematsu2017LC, Poe2020} & medido \\
3 & Modo de fundo: disparos regulares, com taxa de frações a alguns por segundo, que muda com a vigília & Ratos, registro de neurônios do locus coeruleus no ciclo sono-vigília: taxa máxima na vigília, menor no sono de ondas lentas, quase zero no sono REM; as mudanças de taxa antecedem a mudança de estado. & \cite{AstonJonesBloom1981} & medido \\
4 & Modo de rajadas: rajada curta em resposta a um evento importante para a tarefa, mais ou menos no instante da decisão & Macacos, tarefa de encontrar um alvo raro: todos os neurônios do locus coeruleus respondem com rajada só ao alvo, $\sim90\ms$ depois dele e $\sim200\ms$ antes da resposta com a mão; com mau desempenho, as rajadas são mais fracas. Numa tarefa de escolha entre duas opções, a rajada está mais presa à resposta que ao estímulo e ocorre tanto em respostas certas quanto erradas. & \cite{AstonJones1994vigilance, Clayton2004LC} & medido \\
5 & O diâmetro da pupila é um ponto de teste impreciso do \nt{NORA}: ele segue o \nt{NORA}, o \nt{ACET} e outras regiões & Macacos: os disparos do locus coeruleus, até disparos isolados de uma célula, preveem a dilatação da pupila; uma ligação parecida, porém mais tardia, existe nos colículos e no córtex cingulado. Camundongos: as flutuações da pupila acompanham a atividade das saídas noradrenérgicas e colinérgicas no córtex. & \cite{Joshi2016, Reimer2016} & medido \\
6 & A noradrenalina suprime a atividade de fundo mais que a evocada; o ganho multiplicativo~\eqref{eq:gain} é um modelo que reproduz esse efeito & Macacos acordados, córtex auditivo, iontoforese de noradrenalina: a atividade de fundo dos neurônios é mais suprimida que a resposta às vocalizações de coespecíficos, e a relação sinal-ruído cresce. A sigmoide multiplicativa~\eqref{eq:gain} foi tirada de um modelo de rede; a multiplicação literal da inclinação não foi medida. & \cite{Foote1975NE, ServanSchreiber1990} & medido; $g$: modelo \\
7 & O ganho eleva $d'$ só contra o ruído depois do amplificador~\eqref{eq:gainsnr} (proposição~\ref{prop:gainnoise}) & Dedução no capítulo; no modelo de rede, o ganho melhora a detecção do sinal. Não há experimento que separe o ruído antes e depois do amplificador e teste essa diferença. & dedução no capítulo; \cite{ServanSchreiber1990} & teoria \\
8 & A fronteira $g\beta=1$ comuta a rede de escolha de ``pondera'' para ``decidiu''~\eqref{eq:wtagain} (proposição~\ref{prop:gainwta}, fig.~\ref{fig:pitchfork}) & Análise linear de dois grupos que se inibem mutuamente; dedução no capítulo. Não há medições diretas desse limiar no cérebro. & dedução no capítulo & teoria \\
9 & O ganho é o inverso da temperatura da escolha~\eqref{eq:softmax} & Com ruído logístico depois do amplificador, a probabilidade de escolha é um softmax com $T=\sigma/g$; dedução no capítulo. & dedução no capítulo & teoria \\
10 & O \nt{NORA} define quanto buscar: atividade de fundo alta, $g$ efetivo pequeno (busca); rajada no instante da decisão, $g$ grande (decisão) & Humanos: antes de uma escolha ao acaso, a pupila basal é mais larga que antes de escolher a melhor opção conhecida; pessoas com pupila mais larga tendem a buscar mais. Ratos: a passagem para o modo de escolha aleatória é controlada pela entrada do \nt{NORA} no córtex cingulado anterior. Que a rajada eleve o $g$ efetivo não foi medido diretamente. & \cite{JepmaNieuwenhuis2011, Tervo2014stochastic, AstonJones2005} & hipótese \\
11 & A recompensa depende de $g$ como um U invertido; o melhor $g$ é menor num mundo que muda depressa (proposição~\ref{prop:explore}, fig.~\ref{fig:invertedU}) & \texttt{models/nora\_explore.py}, $60$ execuções de $4000$ tentativas. Mudança a cada $1000$ tentativas: melhor $g=16$, fração $0.976$; a cada $20$: $g=8$, fração $0.886$; com $g=0.5$, $0.77$. O recálculo coincidiu com o capítulo. & \texttt{models/\allowbreak nora\_\allowbreak explore.py} & modelo \\
12 & U invertido no comportamento: melhor com atividade de fundo moderada e rajadas nítidas & Macacos, mesma tarefa da linha~4: nos períodos de bom desempenho, a taxa de fundo é moderada e as rajadas ao alvo são nítidas; com taxa de fundo alta não há rajadas e há mais respostas falsas. As mudanças da atividade de fundo e evocada acompanham as oscilações do desempenho. & \cite{Usher1999LC, AstonJones2005} & medido \\
13 & O \nt{NORA} informa a incerteza inesperada; o \nt{ACET}, a esperada & Teoria de inferência bayesiana com dois tipos de incerteza; nos trabalhos citados não há experimento direto que separe esses sinais por neurotransmissor. & \cite{YuDayan2005} & hipótese \\
14 & Depois de uma mudança brusca, as pessoas aprendem mais depressa, e a pupila se dilata & Humanos, tarefa de predição com mudanças súbitas: variações curtas da pupila acompanham a probabilidade estimada de mudança e, junto com a pupila basal, preveem o quanto os novos dados mudam a estimativa (a taxa de aprendizado); uma variação da pupila não ligada à luz nem à tarefa muda essa taxa. Que aqui a pupila mostre justamente o \nt{NORA} é uma inferência a partir dos dados indiretos da linha~5. & \cite{Nassar2012} & medido \\
15 & A rajada funciona como interrupção: a rede zera seu estado & Não há experimento direto em que uma rajada do \nt{NORA} zere o estado da rede; só foram medidas rajadas em eventos importantes (linha~4). & \cite{BouretSara2005} & hipótese \\
16 & O ajuste de $g$ ou da taxa de aprendizado pela surpresa quase não supera as melhores configurações constantes & \texttt{models/nora\_explore.py}, mundo com épocas de $1000$ tentativas (mudança a cada $1000$ e a cada $20$). Melhor constante $g=8$: $0.925$; ajuste de $g$: até $0.927$; ajuste da taxa de aprendizado: até $0.926$; taxa constante $0.4$: $0.928$. O recálculo coincidiu com o capítulo. & \texttt{models/\allowbreak nora\_\allowbreak explore.py} & modelo \\
\bottomrule
\end{longtable}}

\section{Capítulo~\ref{sec:acet}. Acrescentamos o \nt{ACET}}

As três ações da acetilcolina foram medidas em fatias e no cérebro vivo, o conflito entre escrita e leitura foi mostrado no modelo do livro, e que o \nt{ACET} sirva de linha de habilitação de escrita e informe a incerteza esperada é uma hipótese com apoio experimental parcial.

{\footnotesize
\setlength{\tabcolsep}{3pt}\renewcommand{\arraystretch}{1.15}
\begin{longtable}{>{\raggedright}p{0.5cm}>{\raggedright}p{4.0cm}>{\raggedright}p{5.6cm}>{\raggedright}p{2.3cm}>{\raggedright\arraybackslash}p{1.7cm}}
\toprule
N.º & Proposição & Experimento e o que foi medido & Fonte & Status \\
\midrule\endfirsthead
\toprule
N.º & Proposição & Experimento e o que foi medido & Fonte & Status \\
\midrule\endhead
1 & Os neurônios \nt{ACET} do cérebro são poucos, estão reunidos em alguns grupos, e as saídas se espalham por todo o córtex: um canal de difusão & Coloração com anticorpos contra a enzima de síntese da acetilcolina e marcação retrógrada no rato: córtex, hipocampo, amígdala, bulbo olfativo e tálamo recebem acetilcolina principalmente de seis grupos compactos de células do prosencéfalo basal e do tronco & \cite{Mesulam1983} & medido \\
2 & O nível de acetilcolina no córtex é alto na vigília e baixo no sono de ondas lentas & Microdiálise no córtex e no hipocampo de gatos em livre movimento, com registro de EEG: a liberação de acetilcolina na vigília tranquila é $\sim100\%$ e na ativa $\sim175\%$ maior que no sono de ondas lentas; no sono REM, no córtex, é como na vigília & \cite{Marrosu1995,Hasselmo1999} & medido \\
3 & O sentido do sinal é definido pelo receptor: a acetilcolina tem receptores-canais rápidos e receptores lentos, que disparam reações na célula (proposição~\ref{prop:receptor}) & Revisão de medições: o efeito da acetilcolina sobre excitabilidade, transmissão e plasticidade depende do local de liberação, do subtipo de receptor (nicotínicos: canais iônicos; muscarínicos: via proteínas G) e da célula destinatária & \cite{Picciotto2012} & medido \\
4 & Primeira ação: enfraquecer as conexões internas, quase sem afetar as entradas externas (fator $1-a$ em~\eqref{eq:acetnet}) & Fatias de córtex olfativo de rato: com $100$\,\textmu M de agonistas da acetilcolina, os potenciais das conexões internas (camada Ib) caem mais de $60\%$, e os da entrada externa (camada Ia), menos de $15\%$, na mesma célula; a supressão é pré-sináptica. Fatias de hipocampo (CA1): $89\%$ contra $40\%$ para as vias interna e de entrada & \cite{HasselmoBower1992,HasselmoSchnell1994} & medido \\
5 & Segunda ação: reforçar a entrada (fator $\frac{1+a}{2}$) & Fatias de neocórtex: receptores nicotínicos reforçam as sinapses da entrada vinda do tálamo e não agem sobre as intracorticais; os muscarínicos enfraquecem ambos os tipos. A acetilcolina aumenta a excitabilidade dos neurônios (revisão) & \cite{Gil1997,Hasselmo2006} & medido \\
6 & Terceira ação: facilitar a mudança dos pesos (taxa $\eta a$) & Rato: a estimulação elétrica do núcleo basal (fonte de acetilcolina para o córtex) pareada com um tom provoca, no animal adulto, uma forte reorganização do mapa do córtex auditivo em favor do som apresentado & \cite{KilgardMerzenich1998,Hasselmo2006} & medido \\
7 & A regra de Hebb para padrões esparsos na segunda equação de~\eqref{eq:acetnet} dá uma memória associativa que completa o padrão a partir de uma parte & Cálculo da capacidade de uma rede com padrões esparsos e regra de covariância: com fração pequena $p$ de neurônios ativos, a rede guarda muito mais padrões que com $p=1/2$ & \cite{Tsodyks1988} & teoria \\
8 & Escrever com $a$ baixo corrompe a memória: a rede grava o lembrado, e não o visto; com $a=0.1$, a corrupção não é menor que com $a=0$ & \texttt{models/\allowbreak acet\_memory.py}: $200$ neurônios, cinco padrões de $20$, um padrão novo compartilha $10$ neurônios com um deles, $10$ execuções. Com $a=0$, o padrão novo não é memorizado, e o antigo arrasta $5.9$ neurônios alheios de $10$; com $a=0.1$, o novo é lembrado ($0.97$), mas os dois se fundem: o novo arrasta $7.3$ neurônios alheios, o antigo $7.9$; a partir de $a=0.2$, menos de um & dedução no capítulo & modelo \\
9 & Proposição~\ref{prop:writeread}: não há nível $a$ em que escrita e leitura sejam igualmente boas (fig.~\ref{fig:acetsim}) & Mesmo script, taxa de aprendizado $\eta a$: o padrão novo é memorizado por inteiro numa apresentação com $a\ge0.9$, em $0.8$ com $a=0.75$, e não é memorizado com $a\le0.5$. Leitura pela metade do padrão: $1.0$ com $a\le0.2$, $0.96$ com $a=0.3$, $0.04$ com $a=0.4$ & dedução no capítulo & modelo \\
10 & No cérebro, um único fio de difusão, o \nt{ACET}, comuta escrita e leitura & A ideia vem de modelos construídos a partir de medições em fatias. O experimento mais direto é em humanos: o bloqueador de receptores muscarínicos escopolamina prejudica a memorização de novos pares de palavras, sobretudo os que se sobrepõem aos antigos, mas não a lembrança de pares já aprendidos. Não há medição direta dos dois modos da rede & \cite{HasselmoAndersonBower1992,Atri2004} & hipótese \\
11 & O filtro~\eqref{eq:kalman} é ótimo; o peso da entrada e a taxa de aprendizado são a mesma grandeza $P/R$; em regime estacionário, $P=\sqrt{qR}$ e a memória é $\sqrt{R/q}$ (proposição~\ref{prop:kalman}) & Não requer experimento: a otimalidade do filtro de Kalman--Bucy é um resultado clássico da teoria de estimação; o regime estacionário é deduzido no capítulo & \cite{KalmanBucy1961} & teoria \\
12 & O \nt{ACET} informa à rede uma grandeza semelhante a $P$, a incerteza esperada & Proposto num modelo probabilístico de atenção. Não há medição direta de que o nível de acetilcolina acompanhe a incerteza da estimativa & \cite{YuDayan2005} & hipótese \\
13 & Parte dos neurônios \nt{ACET} responde à recompensa e à punição em duas dezenas de milissegundos, tanto mais forte quanto mais inesperado o reforço & Camundongo, neurônios colinérgicos do prosencéfalo basal identificados optogeneticamente: resposta à recompensa e à punição após $18\pm3$\,ms, cuja magnitude cresce com a surpresa do reforço; as respostas são semelhantes nos neurônios de dois núcleos & \cite{Hangya2015} & medido \\
14 & O \nt{NORA} percebe uma mudança inesperada do mundo; o \nt{ACET} mantém a rede no modo de escrita & A divisão de trabalho foi proposta num modelo. Indiretamente: em humanos, o diâmetro da pupila, ligado ao \nt{NORA}, acompanha as mudanças e a taxa de aprendizado. Não há teste direto do par \nt{NORA}--\nt{ACET} & \cite{YuDayan2005,Nassar2012} & hipótese \\
15 & No sono de ondas lentas, a rede em modo de leitura reproduz o que foi gravado de dia, e essas repetições o transcrevem para a memória de longo prazo & Registros de conjuntos de neurônios do hipocampo de rato: células que funcionaram juntas de dia disparam juntas com mais frequência no sono de ondas lentas seguinte, e na mesma ordem: medido. Quanto à transcrição: suprimir as ondulações (ripples) do hipocampo após o aprendizado prejudica a memória espacial, e prolongar as ondulações espontâneas a melhora; que a causa seja o nível baixo de \nt{ACET} não foi demonstrado & \cite{WilsonMcNaughton1994,Skaggs1996,Girardeau2009,FernandezRuiz2019,Hasselmo1999} & hipótese \\
\bottomrule
\end{longtable}}

\section{Capítulo~\ref{sec:synthesis}. A máquina inteira}

O capítulo de síntese não traz experimentos novos: cada linha dele se apoia no capítulo em que a afirmação foi analisada e nas mesmas fontes; o barramento \nt{GLUT} e o controle de fluxo por \nt{GABA} estão firmemente estabelecidos, enquanto os papéis dos canais de difusão e o seu funcionamento conjunto são, em grande parte, hipótese.

{\footnotesize
\setlength{\tabcolsep}{3pt}\renewcommand{\arraystretch}{1.15}
\begin{longtable}{>{\raggedright}p{0.5cm}>{\raggedright}p{4.0cm}>{\raggedright}p{5.6cm}>{\raggedright}p{2.3cm}>{\raggedright\arraybackslash}p{1.7cm}}
\toprule
N.º & Afirmação & Experimento e o que foi medido & Fonte & Status \\
\midrule\endfirsthead
\toprule
N.º & Afirmação & Experimento e o que foi medido & Fonte & Status \\
\midrule\endhead
1 & Para $8.6\cdot10^{10}$ neurônios há apenas quatro tipos de sinais de controle~\eqref{eq:machine} & Contagem de núcleos celulares num homogeneizado de cérebro (fracionador isotrópico): um homem adulto tem $86.1\pm8.1$ bilhões de neurônios & \cite{Azevedo2009} & medido \\
2 & Proposição~\ref{prop:architecture}, as duas primeiras camadas: os dados correm pelo barramento de endereços \nt{GLUT}, o sinal da conexão é definido pelo emissor, o fluxo é dirigido e normalizado pelos neurônios \nt{GABA} (capítulos~\ref{sec:neurontypes}, \ref{sec:gaba}) & Um neurotransmissor principal por neurônio (revisão de medições); a inibição direta estreita a janela em que o neurônio piramidal soma as entradas; a divisão pela atividade total foi medida em muitas regiões & \cite{Strata1999,Pouille2001,Carandini2012} & medido \\
3 & Os elementos não são confiáveis: a sinapse perde a maior parte dos disparos (tab.~\ref{tab:computer}, capítulo~\ref{sec:code}) & Fatias de hipocampo, estimulação mínima: mais da metade dos disparos não provoca resposta em CA1; falhas de limiar e de condução no axônio foram excluídas, a causa é a liberação probabilística do neurotransmissor & \cite{AllenStevens1994} & medido \\
4 & O cérebro funciona com $\sim20$ watts, em média cerca de um disparo por segundo por neurônio (capítulo~\ref{sec:code}); os engenheiros ainda não construíram uma máquina assim & Estimativa~\eqref{eq:energy} a partir dos custos medidos: cerca de $10^9$ moléculas de ATP por disparo, incluindo o trabalho das sinapses (córtex de roedores); uma estimativa detalhada para o córtex dá uma taxa ainda menor. O estado da técnica segundo uma revisão & \cite{Attwell2001,Lennie2003,MehonicKenyon2022} & teoria \\
5 & O aprendizado da maioria das sinapses é o produto de um traço local por um único número comum $\delta$ (segunda equação de~\eqref{eq:machine}, capítulo~\ref{sec:dofa}) & Os neurônios \nt{DOPA} do macaco respondem ao erro de predição de recompensa; em fatias de estriado, a dopamina aumenta a espinha só se chegar $0.3$--$2$\,s depois da entrada glutamatérgica: o traço foi medido & \cite{Schultz1997,Yagishita2014} & medido \\
6 & Um fio de erro próprio para cada saída só existe no circuito de aprendizado motor (capítulo~\ref{sec:motorlearning}) & Cerebelo: a coincidência do sinal da fibra trepadeira com uma entrada enfraquece essa entrada; no macaco, os disparos complexos das células de Purkinje carregam a direção do erro de movimento e provocam a correção & \cite{Ito2001,Herzfeld2018} & medido \\
7 & Cada canal de difusão é um feixe de algumas linhas (proposição~\ref{prop:bundle}) & Para \nt{DOPA}: na mesma tarefa, neurônios diferentes carregam recompensa, movimento e atributos do estímulo; o erro rápido e o nível lento de dopamina são distintos. Para \nt{SERO}, \nt{NORA}, \nt{ACET} essa divisão foi menos estudada & \cite{Engelhard2019,Hamid2016,Mohebi2019} & medido \\
8 & \nt{SERO}: regulador de nível e, por hipótese, o horizonte $\tau_\gamma$ (capítulo~\ref{sec:serocomputer}) & Autoinibição: agonistas dos próprios receptores 5-HT$_{1A}$ inibem os neurônios serotoninérgicos da rafe (medido). O horizonte foi proposto em teoria; o experimento mais forte: a ativação optogenética desses neurônios no camundongo prolonga a espera por uma recompensa adiada & \cite{Sprouse1987,Doya2002,Miyazaki2014} & hipótese \\
9 & \nt{NORA}: o ganho $g$ escolhe entre explorar e decidir (capítulo~\ref{sec:nora}) & Aumento da relação sinal-ruído: no macaco acordado, a noradrenalina no córtex auditivo suprime a atividade de fundo mais do que a resposta a vocalizações de coespecíficos (medido); o ganho multiplicativo é um modelo; o papel na escolha entre explorar e decidir é teoria & \cite{Foote1975NE,ServanSchreiber1990,AstonJones2005} & hipótese \\
10 & \nt{ACET}: alterna entre escrita e leitura (capítulo~\ref{sec:acet}) & Três ações (enfraquecimento das conexões internas, reforço da entrada, facilitação da plasticidade) foram medidas; a troca de modos tem apoio indireto num experimento com escopolamina em humanos & \cite{HasselmoBower1992,Gil1997,Atri2004} & hipótese \\
11 & A fórmula~\eqref{eq:machine} descreve a máquina inteira & É um resumo dos modelos dos capítulos~\ref{sec:glutcomputer}--\ref{sec:acet}, cada parte se apoia no seu capítulo; como um todo, não foi testada experimentalmente & dedução no capítulo & hipótese \\
12 & Os botões trabalham de forma coordenada sem controle comum: erro, surpresa, escrita, retorno (fig.~\ref{fig:timeline}) & Não há experimento: o esquema é qualitativo. Elos isolados: a teoria da divisão de trabalho entre \nt{NORA} e \nt{ACET} e a relação entre pupila e taxa de aprendizado em humanos & \cite{YuDayan2005,Nassar2012} & hipótese \\
13 & O aprendizado por um único sinal comum é tanto mais lento quanto mais neurônios influem no resultado (proposição~\ref{prop:broadcastcost}) & Cálculo das curvas de aprendizado de uma rede linear: o aprendizado por perturbação das saídas é mais lento que a descida de gradiente tantas vezes quantas são as saídas & \cite{Werfel2005} & teoria \\
14 & Não se sabe como o córtex ajusta circuitos de muitas camadas; candidatos: rajadas de disparos, cálculos nos ramos do neurônio, autoaprendizado por predição & Rajadas como portadoras do erro: modelo; reproduzir a resposta de um modelo detalhado de neurônio cortical só uma rede de $5$--$8$ camadas conseguiu: modelo; disparos de cálcio nos ramos de neurônios do córtex humano, que dão o OU exclusivo: medido em fatias; autoaprendizado: revisão de evidências & \cite{Payeur2021,Beniaguev2021,Gidon2020,Keller2018} & hipótese \\
\bottomrule
\end{longtable}}

\section{Capítulo~\ref{sec:sensors}. As entradas da máquina}

A construção dos sensores foi medida diretamente, com registros da corrente de células isoladas, das vibrações da cóclea e com contagens de células da retina; o limite de sensibilidade e a fórmula do ruído balístico decorrem da física, e que os códigos dos sensores estejam ajustados para a máxima informação é uma hipótese com apoio forte.

{\footnotesize
\setlength{\tabcolsep}{3pt}\renewcommand{\arraystretch}{1.15}
\begin{longtable}{>{\raggedright}p{0.5cm}>{\raggedright}p{4.0cm}>{\raggedright}p{5.6cm}>{\raggedright}p{2.3cm}>{\raggedright\arraybackslash}p{1.7cm}}
\toprule
N.º & Proposição & Experimento e o que foi medido & Fonte & Status \\
\midrule\endfirsthead
\toprule
N.º & Proposição & Experimento e o que foi medido & Fonte & Status \\
\midrule\endhead
1 & O sensor é um transdutor: o receptor gera corrente, e a saída das células sensoriais do olho e do ouvido é glutamato para o barramento \nt{GLUT}; as primeiras células não emitem disparos (fig.~\ref{fig:transducer}) & Bastonetes e cones de tartaruga liberam um aminoácido excitatório: ele é captado por canais de glutamato num fragmento de membrana excisado. Células ciliadas internas da cóclea de rato: as correntes nas terminações das fibras nervosas passam por receptores AMPA de glutamato, e a taxa de liberação cresce com a despolarização gradual da célula até $150$\,Hz & \cite{CopenhagenJahr1989,GlowatzkiFuchs2002} & medido \\
2 & O sensor de luz no escuro responde a um único fóton com uma corrente de cerca de um picoampere & Eletrodo de sucção no segmento externo de bastonete de sapo: as respostas a flashes fracos são quantizadas, evento de $\approx1$\,pA ($5\%$ da corrente de escuro) com dispersão de $0.2$\,pA, eficiência $0.5$. Humanos relatam a chegada de um único fóton acima do acaso & \cite{Baylor1979,Tinsley2016} & medido \\
3 & O sensor de som percebe um deslocamento menor que um nanômetro & Método Mössbauer, membrana basilar da cóclea de porquinho-da-índia no ponto de $16$--$19$\,kHz: no limiar de resposta do nervo auditivo, a velocidade da membrana é de cerca de $0.04$\,mm/s, isto é, um deslocamento de $\approx0.35$\,nm (cálculo nosso). Mediu-se a membrana, e não o feixe do sensor & \cite{Sellick1982,Hudspeth1989} & medido \\
4 & No escuro, a precisão é limitada pelo ruído balístico dos fótons~\eqref{eq:photonnoise}; com luz fraca, a acumulação é mais longa & A fórmula decorre da estatística de Poisson. No bastonete, o ruído com luz fraca coincide com a soma de eventos quânticos independentes; sobre um fundo claro, a resposta ao flash é menor e mais curta. O valor ``décimos de segundo'' não foi verificado separadamente aqui & \cite{Baylor1979} & teoria \\
5 & A faixa de entrada é de dez ordens de grandeza para a luz e doze para o som, e a taxa de disparos, menos de três ordens & Para o som: a resposta da membrana basilar na frequência característica cresce com inclinação de até $0.2$\,dB/dB na faixa de $40$--$80$\,dB, e a compressão aparece também acima de $100$\,dB. Os próprios números de ordens são estimativas padrão, sem fonte no capítulo & \cite{Ruggero1997} & medido \\
6 & A resposta do sensor segue~\eqref{eq:nakarushton}, e $\sigma$ acompanha o brilho médio, deslocando a curva ao longo do eixo logarítmico (fig.~\ref{fig:adaptation}) & A fórmula foi ajustada primeiro aos potenciais S da retina de peixe. Cones de tartaruga: as respostas seguem $V=I V_m/(I+\sigma)$ e cobrem $\sim3.5$ ordens de brilho; após $2$--$3$ minutos de fundo, a curva se desloca na horizontal, mantendo a largura. A ligação com a divisão pela atividade total é de uma revisão & \cite{NakaRushton1966,NormannPerlman1979,Carandini2012} & medido \\
7 & Os sensores transmitem mudanças, e os seus códigos estão ajustados para a máxima informação por disparo (proposição~\ref{prop:sensors}) & O princípio foi proposto teoricamente. A evidência mais forte: na mosca, a característica ``contraste $\to$ resposta'' dos neurônios da primeira camada coincide com a distribuição acumulada de contrastes das cenas naturais, e assim se obtém a máxima informação & \cite{Barlow1961,Laughlin1981} & hipótese \\
8 & Os sensores de ligar e de desligar formam um código de dois fios; a câmera de eventos o repete em silício & Revisão de experimentos: os canais de ligar e de desligar da retina se separam já na primeira sinapse e podem ser silenciados separadamente. Câmera: $128\times128$ pixels sem quadros, faixa de $120$\,dB, latência de $15$\,\textmu s & \cite{Schiller1992,Lichtsteiner2008,Gallego2022} & medido \\
9 & O olho tem $\sim10^8$ sensores de luz e $\sim10^6$ neurônios de saída: compressão de $\sim100$ vezes & Contagem em preparações inteiras de retina humana: em média $92$ milhões de bastonetes e $4.6$ milhões de cones; células de saída (ganglionares), de $0.7$ a $1.5$ milhão conforme o olho & \cite{Curcio1990,CurcioAllen1990} & medido \\
10 & O camundongo tem mais de trinta canais de saída do olho & Camundongo: imageamento de cálcio de dois fótons de mais de $11\,000$ células de saída e agrupamento das respostas: bem mais de $30$ tipos funcionais. Para humanos, o número não foi medido & \cite{Baden2016} & medido \\
11 & O olho do porquinho-da-índia transmite $\sim875\,000$ bit/s; o recálculo para o olho humano dá $\sim10^7$ bit/s & Porquinho-da-índia, matriz de múltiplos eletrodos, movimento em cenas naturais: $6$--$13$ bit/s por célula, $\sim875\,000$ bit/s para $\sim10^5$ células de saída. Para humanos ($\sim10^6$ células), $\sim10^7$ bit/s é um recálculo dos autores & \cite{Koch2006} & medido \\
12 & O ouvido é um banco de filtros passa-faixa com um mapa de frequências quase logarítmico (fig.~\ref{fig:filterbank}) & O ponto de máxima vibração da membrana basilar depende da frequência; a função ``frequência--lugar'' é quase exponencial e descreve com uma só forma as cócleas humanas e de animais vivos & \cite{Greenwood1990,Robles2001} & medido \\
13 & Os ressonadores são ativos: parte dos sensores amplifica as vibrações fracas milhares de vezes em amplitude ($66$--$76$\,dB), e o ganho cai com o volume & Vibrometria a laser da base da cóclea de chinchila: com tons fracos na frequência característica, ganho de $66$--$76$\,dB em relação ao estribo; a morte reduz a sensibilidade em $60$--$81$\,dB ($10^3$--$10^4$ vezes em amplitude) e elimina a compressão. A fonte da força são as células ciliadas externas & \cite{Ruggero1997,Robles2001} & medido \\
14 & O amplificador do ouvido está no limite da autoexcitação & Cálculo: perto de um ponto de bifurcação de Hopf são inevitáveis a compressão da faixa, a sintonia aguda e os tons de combinação, isto é, todas as não linearidades conhecidas da audição. Que a cóclea esteja ajustada assim não foi provado diretamente & \cite{Eguiluz2000} & hipótese \\
15 & Nas frequências baixas, os disparos vêm em fase com o som (no porquinho-da-índia, o travamento enfraquece a partir de $\sim600$\,Hz e é perceptível até $\sim3.5$\,kHz), e por eles se acha a direção; nas altas, resta o código de lugar & Nervo auditivo de porquinho-da-índia: o travamento de fase começa a cair em torno de $600$\,Hz e desaparece acima de $3.5$\,kHz; no gato e no macaco o limite fica cerca de uma oitava acima. Diferença de tempo de chegada aos dois ouvidos: revisão & \cite{PalmerRussell1986,Grothe2010} & medido \\
16 & Os demais sensores (tab.~\ref{tab:sensors}): o olfato tem $\sim400$ tipos de receptores, um por neurônio, e código combinatório; o paladar funciona em parte por ATP e serotonina; o tato tem sensores de adaptação rápida e lenta; calor e frio são separados & Camundongo: um receptor reconhece muitos cheiros, um cheiro, muitos receptores, e cheiros diferentes, conjuntos diferentes. Sem receptores de ATP, o nervo do paladar não responde; os botões gustativos liberam serotonina. Registro de fibras isoladas da pele da mão humana: quatro tipos pela velocidade de adaptação. O canal do frio é distinto dos canais do calor & \cite{Mombaerts2004,Malnic1999,Finger2005,Huang2005,JohanssonVallbo1979,McKemy2002} & medido \\
\bottomrule
\end{longtable}}

\section{Capítulo~\ref{sec:recognition}. Reconhecimento}

A velocidade do reconhecimento, a construção dos primeiros estágios, a leitura linear da resposta e o aprendizado pela continuidade do tempo foram medidos em humanos e macacos; a redução de toda a cadeia a duas operações e o aprendizado com vídeo foram testados nos modelos do livro, e as explicações das ilusões vão das bem fundamentadas às controversas.

{\footnotesize
\setlength{\tabcolsep}{3pt}\renewcommand{\arraystretch}{1.15}
\begin{longtable}{>{\raggedright}p{0.5cm}>{\raggedright}p{4.0cm}>{\raggedright}p{5.6cm}>{\raggedright}p{2.3cm}>{\raggedright\arraybackslash}p{1.7cm}}
\toprule
N.º & Afirmação & Experimento e o que foi medido & Fonte & Status \\
\midrule\endfirsthead
\toprule
N.º & Afirmação & Experimento e o que foi medido & Fonte & Status \\
\midrule\endhead
1 & A comparação pixel a pixel com um modelo, sob deslocamento, não acerta mais que uma moeda; o par de estágios~\eqref{eq:scstage} reconhece a figura em qualquer lugar & \texttt{models/\allowbreak recognition\_models.py}, ``retina'' de $24$ pontos, $4000$ tentativas: modelo, $0.49$, $0.51$, $0.52$ com fração de pontos invertidos $0$, $0.1$, $0.2$; par $S$ e $C$, $1.00$, $0.86$, $0.70$ & dedução no capítulo & modelo \\
2 & Um animal numa imagem é reconhecido em $\sim150$\,ms, numa passagem por uma dezena de estágios de $10$--$20$\,ms & Humanos, tarefa ``há um animal ou não'' em fotos mostradas por $20$\,ms: os potenciais evocados das duas respostas divergem após $\sim150$\,ms. Macaco: a latência da primeira resposta cresce de estágio em estágio (V1 primeiro, depois V2 e V4). O número de estágios e ``zero ou um disparo por estágio'' são deduções desses números & \cite{Thorpe1996,Schmolesky1998} & medido \\
3 & O estágio $S$ é um E sobre as partes, o estágio $C$, o máximo sobre as posições (proposição~\ref{prop:conveyor}) & Gato, córtex visual primário: as células simples respondem a uma borda de certa orientação num certo lugar, as complexas, à mesma orientação numa área maior. Registro intracelular de células complexas: a resposta a duas barras, em muitas células, é próxima da maior das duas respostas, em média mais próxima da operação max. Máximo por inibição mútua: proposição~\ref{prop:wta}; divisão pela atividade total: revisão & \cite{HubelWiesel1962,Lampl2004,Carandini2012} & medido \\
4 & Mais acima na cadeia, as combinações são maiores e mais complexas, e a resposta depende cada vez menos da posição & Macaco, simplificação dos estímulos até o ``atributo crítico'': as células de V4 e do córtex inferotemporal posterior respondem a combinações complexas de forma, cor e textura com campo pequeno; no córtex inferotemporal anterior o campo é maior. A alternância de $S$ e $C$ no modelo reproduz esse quadro & \cite{KobatakeTanaka1994,Riesenhuber1999} & medido \\
5 & As redes convolucionais repetem essa alternância e são as que melhor predizem as respostas dos estágios superiores; o objeto é lido linearmente das taxas de um grupo de neurônios & Rede treinada para reconhecer, sem ajuste ao cérebro: a camada superior prediz as respostas dos neurônios do córtex inferotemporal do macaco, as intermediárias, as de V4. Um classificador linear sobre a atividade de $\sim100$ células escolhidas ao acaso, em janelas a partir de $12.5$\,ms, reconhece o objeto e a categoria com posição e tamanho diferentes & \cite{Fukushima1980,Yamins2014,Hung2005} & medido \\
6 & A regra de Hebb com traço~\eqref{eq:traceinvariance} ensina a invariância com vídeo sem rótulos, se o traço não for mais curto que $\sim20$\,ms & A ideia dos atributos lentos e da regra com traço é teoria. \texttt{models/\allowbreak recognition\_models.py}, dois neurônios $C$, $16$ entradas, $10$ rodadas: pausa entre objetos de $0.3$\,s. Com $\tau_{\text{tr}}=5$\,ms, coincidência com a figura de $0.52$ em média, com $10$\,ms, $0.56$; com $20$, $50$ e $200$\,ms, $1.00$ em todas as rodadas (com $30$\,ms, uma rodada em dez erra) & \cite{Foldiak1991,WiskottSejnowski2002} & modelo \\
7 & No cérebro, a invariância é aprendida pela continuidade do tempo & Macaco: o objeto era trocado discretamente durante um salto do olho para um certo lugar do campo; a invariância de posição dos neurônios do córtex inferotemporal mudava de acordo com a troca, de forma significativa já após uma hora. Que essa seja a regra principal do aprendizado visual é hipótese & \cite{LiDiCarlo2008} & medido \\
8 & Movimento: detectores de direção, depois o mesmo par E/OU sobre áreas grandes & Detectores de direção na retina do coelho; na área MT do macaco, todos os $110$ neurônios registrados são seletivos à direção, e a resposta ao movimento é mais forte que em V1 & \cite{Barlow1965,Albright1984} & medido \\
9 & Os estágios superiores enviam para baixo uma predição, e para cima sobe o erro & O modelo explica os efeitos extraclássicos dos campos do córtex visual primário; observações isoladas concordam (revisão); como construção geral da linha de montagem, não está provado & \cite{RaoBallard1999,Keller2018} & hipótese \\
10 & Imagens difíceis exigem realimentação e várias voltas & Macaco: para imagens ``difíceis'', a decisão no córtex inferotemporal se forma $\sim30$\,ms mais tarde; essas respostas tardias são mal preditas por uma rede direta e melhor por redes com realimentação ou mais profundas & \cite{Kar2019} & medido \\
11 & Uma falsificação imperceptível de pixels engana a rede; a visão humana é muito mais robusta, mas não invulnerável & Em redes: uma pequena perturbação escolhida de propósito muda a resposta; o exemplo ``panda $\to$ gibão'' é do segundo trabalho. Em humanos, com apresentação breve, falsificações que se transferem bem entre redes desviam a escolha & \cite{Szegedy2014,Goodfellow2015,Elsayed2018} & medido \\
12 & Ilusões de expectativa: a entrada pouco confiável cede à expectativa; o pálido se move ``mais devagar'', o vestido é visto de formas diferentes (seção~\ref{sec:illusions}) & Grades de menor contraste parecem se mover mais devagar. Enquete com $1401$ pessoas: o vestido é visto em três variantes estáveis, e uma dica sobre a iluminação troca a cor; em $\sim13\,000$ observadores, o que decide é a suposição sobre a iluminação. Um modelo bayesiano com expectativa de movimentos lentos explica várias ilusões de movimento & \cite{Thompson1982,LaferSousa2015,Wallisch2017,Weiss2002,Purves2011} & hipótese \\
13 & Controle de ganho: cachoeira, inclinação e contraste; a grade de Hermann não é inibição dos vizinhos & Os detectores de direção da retina do coelho, com movimento prolongado, reduzem a taxa e, depois da parada, caem abaixo do fundo. A ilusão de inclinação é reproduzida por um modelo com divisão pela atividade dos vizinhos. Mudanças na grade que não alteram a resposta dos neurônios de saída da retina eliminam as manchas: a explicação pela inibição dos vizinhos foi rejeitada & \cite{BarlowHill1963,Schwartz2009,SchillerCarvey2005} & medido \\
14 & Compensação de atrasos: um objeto em movimento parece à frente do clarão & A ilusão foi medida. A psicofísica contradiz tanto o prolongamento do movimento quanto a diferença de latências; foi proposta uma explicação retroativa, pelos eventos $\sim80$\,ms depois do clarão. A disputa não está resolvida & \cite{Nijhawan1994,EaglemanSejnowski2000} & hipótese \\
15 & ``Serpentes'' paradas giram: a diferença de latências~\eqref{eq:latency} é lida pelos detectores de direção & A ilusão funciona também sem cor; pares de elementos vizinhos a geram isoladamente; neurônios do córtex do macaco seletivos à direção dão resposta direcional a esses pares parados. Em humanos, a rotação é disparada por microssacadas e piscadas & \cite{Conway2005,OteroMillan2012} & medido \\
16 & Imagens ambíguas: inibição mútua, fadiga e ruído; as trocas são causadas principalmente pelo ruído & Modelo de grupos de neurônios em competição: nele, as trocas são causadas pelo \emph{ruído}, e sem ruído não existem; ele explica o aumento da frequência das trocas com a força do estímulo. A fadiga tem aqui um papel menor & \cite{MorenoBote2007} & hipótese \\
17 & Parte das ilusões é consequência da tarefa que a máquina aprende & Uma rede treinada para prever o próximo quadro de vídeo ``vê'' rotação em anéis parados e não a vê nas imagens de controle; redes para limpar imagens e aumentar a nitidez caem em ilusões de cor e brilho, como os humanos & \cite{Watanabe2018,GomezVilla2020} & medido \\
\bottomrule
\end{longtable}}

\section{Capítulo~\ref{sec:muscles}. A saída para os músculos}

A construção da saída (unidades motoras, margem de segurança da sinapse, ativação por tamanho, par de músculos, laço de estiramento e geradores de ritmo) foi medida diretamente; a condição de estabilidade do laço com atraso é um teorema; o modelo interno do corpo é uma hipótese bem fundamentada; os números das figuras são cálculos pelo modelo do livro.

{\footnotesize
\setlength{\tabcolsep}{3pt}\renewcommand{\arraystretch}{1.15}
\begin{longtable}{>{\raggedright}p{0.5cm}>{\raggedright}p{4.0cm}>{\raggedright}p{5.6cm}>{\raggedright}p{2.3cm}>{\raggedright\arraybackslash}p{1.7cm}}
\toprule
N.º & Proposição & Experimento e o que foi medido & Fonte & Status \\
\midrule\endfirsthead
\toprule
N.º & Proposição & Experimento e o que foi medido & Fonte & Status \\
\midrule\endhead
1 & Unidade motora: um neurônio \nt{ACET} controla um grupo de fibras, de algumas centenas a um par de milhares; nos músculos de ajuste fino as unidades são menores & Músculos humanos, contagem de axônios motores e de fibras musculares: em média cerca de $340$ fibras por neurônio no músculo interósseo da mão, cerca de $410$ no braquiorradial, cerca de $1930$ no gastrocnêmio medial. O capítulo não dá um número exato para as menores unidades. & \cite{Feinstein1955mu} & medido \\
2 & A sinapse no músculo libera várias vezes mais neurotransmissor do que o necessário para a fibra responder & Revisão de medições em sinapses neuromusculares: a margem de segurança (razão entre o neurotransmissor liberado e o limiar) em mamíferos adultos costuma ser $3$--$5$; em humanos, a margem se apoia mais nas dobras da membrana receptora do que no tamanho da liberação. & \cite{WoodSlater2001} & medido \\
3 & O abalo~\eqref{eq:twitch} sobe num tempo $T_c$ da ordem de $50\ms$ & Unidades motoras isoladas da mão humana em esforço voluntário, força promediada pelos disparos da unidade: tempo de subida de $30$--$100\ms$, abaixo de $70\ms$ em $80\,\%$ das unidades. A função $(t/T_c)e^{1-t/T_c}$ é uma aproximação de um modelo de conjunto de unidades. & \cite{MilnerBrown1973rec, Fuglevand1993} & medido \\
4 & Soma de abalos~\eqref{eq:forcesum}: com $5$, $15$ e $40$ disparos/s a força é $0.2$--$1.1F_1$, $1.8$--$2.2F_1$ e cerca de $5.4F_1$ (fig.~\ref{fig:tetanus}) & Cálculo com \texttt{models/\allowbreak muscle\_\allowbreak models.py}, $T_c=50\ms$: $0.21$--$1.10$, $1.76$--$2.19$ e $5.28$--$5.49\,F_1$ nos últimos $0.3\,\text{s}$. A soma linear é uma simplificação: o modelo não tem a saturação do músculo real. & \texttt{models/\allowbreak muscle\_\allowbreak models.py} & modelo \\
5 & As unidades são ativadas por tamanho; as mais fracas são cem vezes mais fracas que as mais fortes & Raízes ventrais do gato: com entrada reflexa crescente, disparam primeiro os neurônios com impulsos pequenos na raiz, isto é, os pequenos. Músculo interósseo da mão humana: força de abalo das unidades de $0.1$ a $10$~g, crescendo quase linearmente com a força em que a unidade é ativada. & \cite{Henneman1957, MilnerBrown1973rec} & medido \\
6 & O passo de força é proporcional à força, $\Delta F/F\approx q-1$~\eqref{eq:sizeprinciple}: ponto flutuante & Deduzido no capítulo a partir da progressão geométrica das forças. Concorda com o experimento: em esforços grandes, o mesmo incremento de força ativa muito menos unidades novas. & dedução no capítulo; \cite{MilnerBrown1973rec} & teoria \\
7 & A dispersão da força cresce proporcionalmente à própria força & Esforço isométrico voluntário do músculo do polegar: o desvio padrão da força cresce linearmente com a força média; com o mesmo esforço provocado por estimulação elétrica, a dispersão é constante. Modelo de conjunto: o crescimento linear resulta da ativação das unidades por ordem de força. & \cite{Jones2002sdn} & medido \\
8 & A suavidade dos movimentos é consequência do ruído dependente do sinal & Modelo: a trajetória que minimiza a dispersão do ponto final com esse ruído reproduz as trajetórias medidas do olho e do braço. Não há experimento direto que distinga essa causa das outras. & \cite{HarrisWolpert1998} & hipótese \\
9 & A diferença das forças do par de músculos dá o torque, a soma dá a rigidez~\eqref{eq:antagonists} & Teoria do controle de impedância por co-contração. Braço humano: a rigidez de cada articulação, medida com pequenos empurrões, é aproximadamente proporcional ao torque nela; diferentes proporções de co-contração mudam a forma e a direção da elipse de rigidez da mão. & \cite{Hogan1984, GomiOsu1998stiff} & medido \\
10 & O sensor de estiramento excita o seu neurônio \nt{ACET} e, por um intermediário inibitório, desliga o antagonista (fig.~\ref{fig:antagonists}) & Medula espinhal do gato: um único intermediário, excitado pelas fibras dos sensores de estiramento, produz potenciais inibitórios monossinápticos nos neurônios \nt{ACET} do músculo antagonista, atraso sináptico de $0.28$--$0.42\ms$. O neurotransmissor do intermediário (glicina) não foi determinado nesse experimento. & \cite{Jankowska1972Ia} & medido \\
11 & Atraso do laço: espinhal $25$--$30\ms$, pelo córtex uma vez e meia a três vezes mais, pelo olho $100$--$200\ms$ & Braço humano, empurrão na articulação: resposta muscular curta após $20$--$45\ms$, longa após $45$--$105\ms$, voluntária após $120$--$180\ms$. Deslocamento da posição visível do dedo durante o movimento: correção após $160\ms$ em média. & \cite{Pruszynski2008rapid, SaundersKnill2003vis} & medido \\
12 & $\dd x/\dd t=-Gx(t-d)$ é estável com $Gd<\pi/2$~\eqref{eq:delaystable}; no limite, a frequência é $1/(4d)$ & Teorema sobre as raízes de uma equação com retardo. Verificação com \texttt{models/\allowbreak muscle\_\allowbreak models.py}, $d=30\ms$: em $3\,\text{s}$, amplitude $3\cdot10^{-6}$ com $G=40$, $0.13$ com $G=50$, $38$ com $G=55$ por segundo (limite $52$). & \cite{Hayes1950} & teoria \\
13 & Tremor fisiológico de $8$--$12$~Hz; o laço de estiramento é uma das suas fontes & Revisão de medições: um tremor fino na faixa de cerca de $10$~Hz existe em pessoas saudáveis; a sua origem é mista (oscilações centrais, propriedades das descargas das unidades, ressonância mecânica e ressonância do laço reflexo), e em condições diferentes predomina uma ou outra. & \cite{McAuleyMarsden2000} & hipótese \\
14 & O cérebro prevê o resultado do comando por um modelo interno do corpo & Uma pessoa segura um peso em pinça e move o braço: com carga inercial, viscosa e elástica, a força de preensão muda ao mesmo tempo que a carga, e não com o atraso do laço, ou seja, a carga é prevista. Uma revisão reúne esses experimentos; não há observação direta do próprio modelo nos neurônios. & \cite{FlanaganWing1997grip, WolpertGhahramani2000} & hipótese \\
15 & O ritmo de andar, respirar e mastigar é gerado por redes locais com entrada constante & Revisão de experimentos: o sistema nervoso isolado (medula, gânglios) produz um padrão motor rítmico sem entrada rítmica e sem realimentação dos músculos. & \cite{MarderBucher2001} & medido \\
16 & A inibição mútua com fadiga~\eqref{eq:halfcentre} dá um ritmo com período de $0.84\,\text{s}\approx2\tau_a$ (fig.~\ref{fig:halfcentre}) & Teorema: uma rede com inibição mútua e adaptação sem estado estável, com entrada constante, necessariamente oscila. Cálculo com \texttt{models/\allowbreak muscle\_\allowbreak models.py} ($I=1$, $\beta=2$, $b=2$, $\tau=20\ms$, $\tau_a=0.4\,\text{s}$): período de $0.84\,\text{s}$. Que os geradores reais sejam construídos exatamente assim não foi mostrado. & \cite{Matsuoka1985osc} & modelo \\
\bottomrule
\end{longtable}}

\section{Capítulo~\ref{sec:pain}. A dor}

Os sensores de dor, os dois fios, a porta na medula espinhal, o botão de volume descendente, a sensibilização e a captura da atenção foram medidos diretamente; as fórmulas da porta e da sensibilização são modelos simplificados do livro; a regra pela qual a máquina escolhe o volume do alarme é hipótese.

{\footnotesize
\setlength{\tabcolsep}{3pt}\renewcommand{\arraystretch}{1.15}
\begin{longtable}{>{\raggedright}p{0.5cm}>{\raggedright}p{4.0cm}>{\raggedright}p{5.6cm}>{\raggedright}p{2.3cm}>{\raggedright\arraybackslash}p{1.7cm}}
\toprule
N.º & Proposição & Experimento e o que foi medido & Fonte & Status \\
\midrule\endfirsthead
\toprule
N.º & Proposição & Experimento e o que foi medido & Fonte & Status \\
\midrule\endhead
1 & Limiar alto: os limiares de calor dos diferentes sensores de dor vão de $41^\circ$ a $46$--$48^\circ$ & Registro de fibras isoladas. Pele de macaco: limiar mediano dos sensores lentos (C) de $41^\circ$, dos sensores rápidos do segundo tipo, $46^\circ$, do primeiro tipo, acima de $53^\circ$. Pele humana: sensores C que respondem também à pressão, $40.7^\circ$; os que não respondem à pressão, $48^\circ$. & \cite{Treede1995heat, Weidner1999cfib} & medido \\
2 & Os sensores de dor se adaptam muito menos que os de toque e ficam mais sensíveis após uma lesão leve; o enfraquecimento após aquecimento forte foi medido num tipo & Queimadura leve da pele ($50^\circ$, $100\,\text{s}$): após $5$--$10$~min, o limiar de dor em humanos e o limiar dos sensores C no macaco caem $2$--$6^\circ$; os outros sensores da pele não mostram esse desvio. Nos sensores rápidos do segundo tipo, a resposta após aquecimento forte fica suprimida. A comparação da adaptação com os sensores de toque é uma estimativa geral, sem experimento próprio aqui. & \cite{LaMotte1982hyper, Treede1995heat} & medido \\
3 & Dois fios: rápido de $5$--$30$~m/s, lento de $0.5$--$2$~m/s; tempo de chegada $t=L/v$~\eqref{eq:paintime} & Velocidade de condução: sensores de dor rápidos do macaco, em média $15$ e $25$~m/s (dois tipos); sensores C humanos, $0.8$--$1.0$~m/s. Com $L=1$~m, isso dá dezenas de milissegundos e cerca de um segundo. & \cite{Treede1995heat, Weidner1999cfib} & medido \\
4 & A dor chega duas vezes: primeiro rápida e precisa, depois surda e longa & Tempo de reação ao aquecimento do antebraço humano: de $1100$ a $400\ms$ com o aumento da temperatura; o aquecimento forte da pele com pelos dá dor dupla, o da palma, não. A primeira dor só pode ser levada por fibras mais rápidas que $6$~m/s: pela latência, correspondem a ela os sensores rápidos do segundo tipo, que não existem na palma. A divisão de papéis (``onde'' e ``quão ruim'') é interpretação. & \cite{CampbellLaMotte1983first, Treede1995heat} & medido \\
5 & Porta: o neurônio de saída da medula recebe dor e toque, e o intermediário inibitório é excitado pelo toque (fig.~\ref{fig:gate}) & O esquema da porta foi proposto como teoria. Camundongos, marcação genética e remoção de grupos de neurônios: os neurônios excitatórios com somatostatina são necessários para a dor mecânica; os neurônios inibitórios com dinorfina impedem que a entrada dos sensores de toque chegue a eles; sem esses neurônios, um toque leve provoca dor. & \cite{MelzackWall1965, Duan2014} & medido \\
6 & O toque abafa a dor & Humanos, pulsos de laser dolorosos na mão: o toque simultâneo reduz a avaliação da dor e a capacidade de distinguir a energia dos pulsos; o efeito enfraquece com a distância entre o ponto do toque e o ponto da dor dentro de um mesmo segmento da pele. & \cite{Mancini2014touch} & medido \\
7 & Suposição da teoria da porta: uma dor forte desliga por si o intermediário inibitório, e a porta se escancara & No capítulo, marcado como suposição da teoria original. O trabalho em camundongos descreve como a porta impede o toque de entrar no canal da dor; o desligamento do intermediário pela dor não é mostrado nele. & \cite{MelzackWall1965} & hipótese \\
8 & Porta~\eqref{eq:gate}: limiar de $0.18$ sem toque, $0.86$ com toque, $0.11$ com $g=2$; com $\nu=1$ o toque reduz a saída de $1$ para $0.25$, com $\nu=2$ não tem efeito; o toque sozinho não passa (fig.~\ref{fig:gatecurves}) & O cálculo com \texttt{models/\allowbreak pain\_\allowbreak gate.py}, com os pesos do texto, reproduz todos esses números. & \texttt{models/\allowbreak pain\_\allowbreak gate.py} & modelo \\
9 & As vias descendentes do tronco mudam o ganho da porta nos dois sentidos & Rato, bulbo, registro durante aquecimento da cauda: alguns neurônios disparam antes da retirada (``on''), outros se calam (``off''); uma estimulação fraca dessa região suprime o reflexo. Revisão: os mesmos circuitos tanto facilitam quanto inibem a transmissão da dor, e os opioides agem por meio deles. & \cite{Fields1983rvm, Fields2004} & medido \\
10 & A expectativa de alívio abafa a dor pelo canal opioide & Pessoas com dor aguda: naquelas em que a expectativa de alívio reduziu a dor, o bloqueio dos receptores opioides devolveu a dor ao nível das demais; nas demais, o bloqueio não mudou a dor. & \cite{Levine1978} & medido \\
11 & O cérebro escolhe o volume do alarme pelo benefício da interrupção & Não há experimento em que a regra de escolha do volume tenha sido medida. & --- & hipótese \\
12 & Sensibilização: após uma lesão, a saída da porta cresce e o limiar cai, em parte por causa da medula; ela persiste depois que o estímulo lesivo cessa & Rato: após lesão da pele, o limiar do reflexo de flexão cai e a resposta cresce; a análise eletrofisiológica mostra que parte do aumento surge na própria medula. O estímulo lesivo provoca alterações duradouras da sensibilidade, que sobrevivem ao próprio estímulo. O capítulo não dá o tempo exato de decaimento. & \cite{Woolf1983} & medido \\
13 & A sensibilização~\eqref{eq:sensitization} é uma realimentação positiva; com um laço forte, o alarme pode travar depois da cura & Decorre do modelo~\eqref{eq:sensitization} (exercício no fim do capítulo). Não há experimento mostrando que a dor persistente após a cura seja um laço travado exatamente desse tipo. & dedução no capítulo & hipótese \\
14 & A dor é uma recompensa negativa, e o aprendizado com ela segue o erro de diferença temporal & Ressonância funcional, humanos, cadeias de prenúncios da dor: a atividade do estriado ventral e da ínsula anterior coincide com os sinais de aprendizado por diferença temporal. Macaco: parte dos neurônios \nt{DOPA} é excitada por um sopro de ar desagradável e pelo seu prenúncio, outra parte é inibida. & \cite{Seymour2004, Matsumoto2009} & medido \\
15 & A dor interrompe o trabalho em curso & Humanos, estímulo elétrico doloroso e testes sonoros: a resposta ao teste $100\ms$ após o início da dor fica bem mais lenta; após $1.5\,\text{s}$ já não há diferença em relação ao controle. Uma revisão reúne esses experimentos num modelo de captura da atenção. & \cite{Crombez1996disrupt, EcclestonCrombez1999} & medido \\
16 & A retirada se fecha na medula espinhal & Rato: neurônios das camadas profundas do corno dorsal reproduzem o padrão espacial do reflexo de retirada de músculos individuais; eles não podem ser excitados antidromicamente a partir da região cervical, ou seja, são intermediários espinhais. & \cite{Schouenborg1995wr} & medido \\
\bottomrule
\end{longtable}}

\section{Capítulo~\ref{sec:motorlearning}. O aprendizado motor}

A regra dos mínimos quadrados e a vantagem de um fio de erro separado são teoremas; a construção do circuito com fio de erro, o enfraquecimento por coincidência, o ajuste do traço ao atraso, o pós-efeito e o retorno espontâneo da habilidade foram medidos; que o circuito inteiro funcione como aprendizado supervisionado e construa um modelo interno é a hipótese principal; os números do modelo de dois processos são cálculos pelo modelo do livro.

{\footnotesize
\setlength{\tabcolsep}{3pt}\renewcommand{\arraystretch}{1.15}
\begin{longtable}{>{\raggedright}p{0.5cm}>{\raggedright}p{4.0cm}>{\raggedright}p{5.6cm}>{\raggedright}p{2.3cm}>{\raggedright\arraybackslash}p{1.7cm}}
\toprule
N.º & Afirmação & Experimento e o que foi medido & Fonte & Status \\
\midrule\endfirsthead
\toprule
N.º & Afirmação & Experimento e o que foi medido & Fonte & Status \\
\midrule\endhead
1 & A regra~\eqref{eq:lms} segue, em média, o gradiente do erro quadrático & Resultado da teoria dos filtros adaptativos: uma correção proporcional à entrada e ao erro da saída é uma descida estocástica pelo gradiente do erro quadrático médio. & \cite{WidrowHoff1960} & teoria \\
2 & Com um fio de erro para cada saída, a velocidade não depende do número de saídas; com um único número comum, cai com o número de neurônios ruidosos (proposição~\ref{prop:teachers}) & Curvas de aprendizado de uma rede linear: o aprendizado por perturbações aleatórias dos neurônios é mais lento que o gradiente exato tantas vezes quantos são os neurônios perturbados. & \cite{Werfel2005} & teoria \\
3 & O circuito com fio de erro funciona como aprendizado supervisionado (proposição~\ref{prop:errorwire}) & Teorias deduzidas da estrutura do circuito. Os experimentos das linhas 7--10 concordam com elas; que o circuito inteiro funcione exatamente assim não está provado. & \cite{Marr1969, Albus1971} & hipótese \\
4 & Camada de expansão: cerca de $7\cdot10^{10}$ pequenos neurônios \nt{GLUT}, mais que em todo o resto do cérebro & Contagem de núcleos numa suspensão de tecido dissolvido: no cerebelo humano há $69\cdot10^9$ neurônios dos $86\cdot10^9$ do cérebro inteiro. Estereologia: $101\cdot10^9$ pequenos neurônios no cerebelo de homens idosos. & \cite{Azevedo2009, Andersen1992cereb} & medido \\
5 & Aumentar a dimensão da entrada torna linear a dependência desejada & Teorema sobre o número de dicotomias: uma soma ponderada de $n$ entradas com limiar realiza uma rotulação arbitrária de cerca de $2n$ vetores em posição geral. Que o cerebelo use exatamente isso é parte da hipótese da linha~3. & \cite{Cover1965} & teoria \\
6 & Cerca de $3\cdot10^7$ neurônios \nt{GABA} de saída, cada um com da ordem de $10^5$ entradas & Estereologia do cerebelo humano: $30.5\cdot10^6$ neurônios de saída. Microscopia eletrônica, rato: cerca de $1.75\cdot10^5$ entradas da camada de expansão por neurônio de saída. O neurotransmissor inibitório dos neurônios de saída está na revisão. & \cite{Andersen1992cereb, NapperHarvey1988, Ito2001} & medido \\
7 & Exatamente um fio de erro por neurônio de saída, cerca de uma vez por segundo, cada vez uma descarga poderosa & Revisão de registros: cada neurônio de saída recebe uma entrada \nt{GLUT} trepadeira, que provoca um disparo complexo com frequência da ordem de $1$~Hz. & \cite{Ito2001} & medido \\
8 & A probabilidade da descarga depende da direção do erro de movimento & Macaco, região oculomotora do cerebelo, adaptação de sacadas: a direção do erro visual define a probabilidade do disparo complexo, e a magnitude, o seu momento; a descarga muda os disparos simples na tentativa seguinte e desloca o movimento ao longo do erro preferido da célula. & \cite{Herzfeld2018} & medido \\
9 & A coincidência da descarga de erro com uma entrada ativa enfraquece a entrada ($-\eta e_i x_j$) & Revisão de experimentos: a estimulação conjunta de uma entrada da camada de expansão e do fio de erro provoca um enfraquecimento de longo prazo dessa entrada; a estimulação de uma só entrada, não. & \cite{Ito2001} & medido \\
10 & O comprimento do traço é ajustado ao atraso do erro: com atraso de $100\ms$, enfraquece mais a entrada que esteve ativa $100\ms$ antes & Fatias e camundongos vivos: na região responsável pelo aprendizado oculomotor, o enfraquecimento é sintonizado estreitamente num intervalo de cerca de $120\ms$ entre entrada e erro, o tempo que o sinal de erro leva nessa tarefa; em outra região, células diferentes são sintonizadas em intervalos diferentes. & \cite{Suvrathan2016} & medido \\
11 & O circuito é um filtro adaptativo~\eqref{eq:adaptivefilter} que aprende o modelo interno do corpo & Modelo deduzido da estrutura do circuito. Não há experimento direto em que o filtro aprendido tenha sido medido como função de transferência do corpo. & \cite{Fujita1982} & hipótese \\
12 & A predição subtrai as consequências dos próprios movimentos, por isso não conseguimos fazer cócegas em nós mesmos & Humanos: o próprio toque parece menos cócegas que o mesmo toque alheio; na ressonância, o córtex somatossensorial responde mais fraco ao próprio toque, e o cerebelo fica menos ativo quando o movimento toca a pele. A explicação pela subtração da predição é hipótese. & \cite{Blakemore1998} & hipótese \\
13 & Com uma força externa, o erro diminui e, após a sua retirada, a mão erra para o outro lado & Pessoas movem a manopla de um robô num campo de forças: as trajetórias voltam a ser retas; depois da retirada do campo, o pós-efeito é quase um espelho das primeiras distorções; ele se transfere também para partes não treinadas do espaço. & \cite{ShadmehrMussaIvaldi1994} & medido \\
14 & Aprendem dois processos~\eqref{eq:tworate}; depois da perturbação inversa, a habilidade volta sozinha & Humanos, campo de forças, depois um campo inverso curto e tentativas com o erro travado: a correção volta sozinha à antiga; o modelo de dois processos reproduz isso, o de um só, não. & \cite{Smith2006} & medido \\
15 & Números da fig.~\ref{fig:tworate}: $0.84$ (lento $0.63$) em $200$ movimentos; $6$ inversos; rápido $-0.53$, lento $0.46$; retorno até $0.37$ em $40$ movimentos & Cálculo com \texttt{models/\allowbreak motor\_\allowbreak learning.py}, $A_f=0.92$, $B_f=0.1$, $A_s=0.996$, $B_s=0.02$: $0.840$ ($0.634$ e $0.205$), $6$ movimentos, $-0.531$ e $0.457$, pico de $0.371$ no $40$º movimento. & \texttt{models/\allowbreak motor\_\allowbreak learning.py} & modelo \\
16 & Dois processos são necessários porque o corpo muda em velocidades diferentes & Teoria: a estimativa ótima com perturbações de vários tempos de vida dá vários filtros com constantes de tempo diferentes e explica o retorno espontâneo e o reaprendizado acelerado. & \cite{Kording2007} & hipótese \\
\bottomrule
\end{longtable}}

\section{Capítulo~\ref{sec:chemoutput}. A saída química}

Os dois fios para o coração e suas velocidades diferentes, a realimentação negativa nas cascatas hormonais, o código pela frequência das porções, o gerador diário e seu ajuste pela luz foram medidos diretamente; o limite $K<8$ é um teorema da teoria de controle, deduzido no capítulo; as pulsações geradas pela própria realimentação da cascata são uma hipótese com um experimento forte a seu favor.

{\footnotesize
\setlength{\tabcolsep}{3pt}\renewcommand{\arraystretch}{1.15}
\begin{longtable}{>{\raggedright}p{0.5cm}>{\raggedright}p{4.0cm}>{\raggedright}p{5.6cm}>{\raggedright}p{2.3cm}>{\raggedright\arraybackslash}p{1.7cm}}
\toprule
N.º & Proposição & Experimento e o que foi medido & Fonte & Status \\
\midrule\endfirsthead
\toprule
N.º & Proposição & Experimento e o que foi medido & Fonte & Status \\
\midrule\endhead
1 & O marca-passo do coração bate sozinho cerca de $100$ vezes por minuto; em repouso o freio está pisado, e a frequência costuma ficar em torno de $60$--$70$ & Seres humanos, desligamento completo dos dois fios para o coração: a frequência própria é maior que a de repouso e diminui com a idade; em adultos, da ordem de $100$ batimentos por minuto. Logo, em repouso predomina o freio. A frequência de repouso de $60$--$70$ é um valor típico, sem citação à parte. & \cite{JoseCollison1970ihr} & medido \\
2 & O acelerador \nt{NORA} e o freio \nt{ACET} são filtros passa-baixas, e o freio é mais rápido~\eqref{eq:gasbrake} & Cão, estimulação modulada em frequência do nervo vago ou do simpático cardíaco e função de transferência até a frequência atrial: as duas vias são filtros passa-baixas; a simpática tem frequência de corte muito menor e um atraso puro de cerca de $1.7\,\text{s}$. As características dependem do tônus médio, de modo que a fórmula linear~\eqref{eq:gasbrake} é uma aproximação. & \cite{Berger1989} & medido \\
3 & O freio é rápido porque o \nt{ACET} abre o canal de potássio sem segundo mensageiro, e o \nt{NORA} age por uma cadeia de reações & Fragmentos de membrana de células atriais: o canal de potássio aberto pela acetilcolina é aberto pelas subunidades $\beta\gamma$ da proteína G ligada ao receptor, sem segundo mensageiro dentro da célula. A via do \nt{NORA} pelo AMPc não é citada aqui. & \cite{Logothetis1987gbg} & medido \\
4 & O sangue dá a volta no corpo em cerca de um minuto; o \nt{ADRE} da corrente sanguínea chega a todos os órgãos com receptor & Estimativa geral de ordem de grandeza; assim está formulada no capítulo, sem fonte citada. & --- & estimativa \\
5 & A cascata hormonal é envolvida por realimentação negativa: o hormônio final inibe os primeiros estágios & Revisão de experimentos: os hormônios do córtex adrenal suprimem a liberação de ACTH pela hipófise e do hormônio liberador pelo hipotálamo em três faixas de tempo: rápida, intermediária e lenta. & \cite{KellerWood1984fb} & medido \\
6 & Três estágios iguais são estáveis para $K<8$~\eqref{eq:cascadestable}; no limite, o período é $2\pi\tau/\sqrt3\approx3.6\tau$ & Deduzido no capítulo: na frequência $\sqrt3/\tau$, os três estágios dão defasagem de $180^\circ$ e atenuação de $8$ vezes. & dedução no capítulo & teoria \\
7 & O erro de nível é $1/(1+K)$, por isso esse regulador não tem precisão melhor que uns $10\,\%$ (proposição~\ref{prop:cascade}) & Consequência da linha~6 para realimentação proporcional. O eixo real tem realimentação em vários estágios e em várias faixas de tempo (linha~5), de modo que o limite se refere ao modelo~\eqref{eq:cascade} e não foi medido. & dedução no capítulo & teoria \\
8 & Com $K=4$ e $7$ o nível se estabiliza; com $K=12$, pulsações regulares com período de $3.7$--$3.9\,\tau$ (fig.~\ref{fig:cascade}) & Cálculo com \texttt{models/\allowbreak chem\_\allowbreak models.py}: com $K=12$, períodos sucessivos de $3.9$, $3.7$, $3.8$, $3.7\,\tau$; com $K=4$, amplitude de $0.003$ no fim do cálculo. & \texttt{models/\allowbreak chem\_\allowbreak models.py} & modelo \\
9 & O hormônio do estresse sai em pulsos cerca de uma vez por hora; os pulsos nascem da própria realimentação, sem gerador separado & Ratos, infusão constante de hormônio liberador: ACTH e corticosterona oscilam com frequência fisiológica, quase horária; uma entrada rítmica do hipotálamo não é necessária para os pulsos. Um modelo hipófise--adrenal com realimentação atrasada produz essas oscilações. & \cite{Walker2012glc, Walker2010} & hipótese \\
10 & O destinatário lê a frequência das porções, não o nível & Macacos sem liberação própria de hormônio liberador de gonadotrofinas: a infusão constante não restabelece a liberação dos hormônios da hipófise, as porções de hora em hora restabelecem; a volta à infusão constante suprime de novo a resposta. A fadiga do receptor como filtro passa-altas é uma interpretação. & \cite{Belchetz1978} & medido \\
11 & Frequências diferentes de porções agem de modo diferente sobre células diferentes da mesma glândula & Os mesmos macacos: aumentar as porções para $2$--$5$ por hora reduz os dois hormônios da hipófise; reduzi-las a uma a cada $3$~horas diminui um (LH) e aumenta o outro (FSH), de modo que a razão entre eles é definida pela frequência. & \cite{Wildt1981gnrh} & medido \\
12 & O ritmo diário nasce de uma realimentação negativa intracelular com atraso de algumas horas & Revisão: as proteínas dos genes do relógio suprimem, com atraso, a transcrição dos próprios genes; o laço de transcrição e tradução dá um período de cerca de um dia. & \cite{TakahashiJS2017} & medido \\
13 & O gerador principal é um grupo de dezenas de milhares de neurônios que sincroniza o resto do corpo & Revisão: o núcleo supraquiasmático é o gerador diário principal; seus neurônios isolados em cultura são geradores autônomos, e a rede os sincroniza; no camundongo, cerca de $10^4$ neurônios de cada lado. & \cite{Welsh2010scn} & medido \\
14 & O período próprio no ser humano é de $24.18$~horas & Pessoas sob luz fraca, com horário desvinculado do dia: o período dos ritmos de melatonina, temperatura e cortisol foi em média de $24.18$~horas em jovens e idosos, com dispersão estreita. & \cite{Czeisler1999} & medido \\
15 & A luz acerta o gerador como uma malha de captura de fase~\eqref{eq:pll}; é preciso um deslocamento de cerca de $11$~min por dia & Pessoas, um único pulso de luz intensa de $6.7$~horas em diferentes horas do dia: curva de resposta de fase do tipo~1 com amplitude de $5$~horas: atraso antes do mínimo da temperatura corporal, avanço depois. A equação~\eqref{eq:pll} é o modelo padrão; $11$~min $=0.18$~hora é aritmética. & \cite{Khalsa2003prc} & medido \\
16 & Sem luz não há ajuste, e o ritmo segue o período próprio & Pessoas totalmente cegas, sem percepção de luz, vivendo em sociedade normal: em cerca de metade delas os ritmos de melatonina e cortisol seguem um período próprio, que não coincide com o dia. & \cite{Sack1992blind} & medido \\
\bottomrule
\end{longtable}}

\section{Capítulo~\ref{sec:debugging}. Depurando a máquina}

O que o eletrodo ouve, as poucas dimensões da atividade, o desempenho das interfaces com matrizes rígidas, o aprendizado conjunto do cérebro e do decodificador e os pontos visuais produzidos por estimulação foram medidos em experimentos revisados por pares; as estimativas dos fluxos de dados são aritmética do capítulo; sobre o dispositivo da Neuralink implantado em pessoas, até onde foi possível verificar, os dados foram publicados sobretudo pela própria empresa.

{\footnotesize
\setlength{\tabcolsep}{3pt}\renewcommand{\arraystretch}{1.15}
\begin{longtable}{>{\raggedright}p{0.5cm}>{\raggedright}p{4.0cm}>{\raggedright}p{5.6cm}>{\raggedright}p{2.3cm}>{\raggedright\arraybackslash}p{1.7cm}}
\toprule
N.º & Proposição & Experimento e o que foi medido & Fonte & Status \\
\midrule\endfirsthead
\toprule
N.º & Proposição & Experimento e o que foi medido & Fonte & Status \\
\midrule\endhead
1 & O eletrodo ouve os disparos dos neurônios num raio da ordem de cem micrômetros, em geral de um a alguns; eles são distinguidos pela forma & Rato, campo CA1, registro intracelular e extracelular simultâneo: a forma do disparo extracelular reproduz a largura e a amplitude do intracelular. Estimativa: um tetrodo poderia distinguir $60$--$100$ neurônios, mas na prática se isolam cerca de seis. & \cite{Henze2000extra} & medido \\
2 & O cérebro tem $8.6\cdot10^{10}$ neurônios; a sonda ouve cerca de um centésimo de milionésimo & Contagem de núcleos numa suspensão de tecido dissolvido: $86\cdot10^9$ neurônios no cérebro humano. A fração é aritmética do capítulo. & \cite{Azevedo2009} & medido \\
3 & A atividade de centenas de neurônios do córtex motor cabe em uma ou duas dezenas de variáveis comuns & Revisão de registros no córtex motor de macacos: a atividade da população é descrita por poucas variáveis latentes, comuns a muitos neurônios. & \cite{Gallego2017} & medido \\
4 & O ruído dos vizinhos é correlacionado, e cem neurônios carregam quase tanto quanto mil (proposição~\ref{prop:pooling}) & Córtex visual de macaco: coeficiente de correlação do ruído de neurônios vizinhos em torno de $0.1$, e o ganho da média satura com uma centena de neurônios. Teoria das correlações que limitam a informação. & \cite{Zohary1994, MorenoBote2014} & medido \\
5 & Fluxo bruto de $2\cdot10^8$~bit/s contra $8\cdot10^4$~bit/s em disparos~\eqref{eq:probedata} & Aritmética do capítulo pela capacidade do fluxo de disparos~\eqref{eq:spikeentropy}. Os parâmetros de digitalização são reais: o chip da Neuralink usa $19.3$~kHz e $10$~bits por canal. & dedução no capítulo; \cite{Rieke1997, Musk2019} & teoria \\
6 & Primeiro sistema da Neuralink: os chips amplificam e digitalizam, o fluxo bruto sai por cabo, os disparos são detectados por um programa externo; detecção no chip, carcaça fechada, rádio e energia sem fio estão no dispositivo para pessoas, segundo a empresa & Artigo da empresa: o fluxo bruto completo sai por cabo USB-C, e os disparos são detectados em tempo real por um programa fora do chip; a compressão embarcada, a alimentação e a transmissão sem fio são citadas como plano para dispositivos clínicos. Para o dispositivo implantado em pessoas, só há comunicados da empresa. Que a detecção de disparos no chip é necessária é uma conclusão do capítulo a partir de~\eqref{eq:probedata}. & \cite{Musk2019}; relatórios da Neuralink, sem artigo revisado por pares & medido (artigo da empresa); em pessoas, relatório da empresa \\
7 & Até $3072$ eletrodos em $96$ fios de $32$; os fios são inseridos por um robô que desvia dos vasos & Artigo da empresa: $3072$ eletrodos em $96$ fios de $32$; o robô insere $6$ fios ($192$ eletrodos) por minuto; em ratos com matriz implantada por longo prazo, até $70\,\%$ dos eletrodos ouvem disparos. & \cite{Musk2019} & medido (artigo da empresa) \\
8 & Em pessoas, cerca de mil canais em $64$ fios; parte dos fios se deslocou; cursor da ordem de $10$~bit/s & Só comunicados da empresa. Uma busca no PubMed não encontrou artigo revisado por pares com resultados em pessoas; isso concorda com a ressalva no texto do capítulo. & relatórios da Neuralink; sem artigo revisado por pares & medido (relatório da empresa) \\
9 & Uma interface com matriz de cem eletrodos controla um cursor & Pessoa com paralisia, $96$ eletrodos no córtex motor primário: a intenção de mover o braço altera os disparos três anos após a lesão; o decodificador fornece um ``cursor neural'' (e-mail, televisão), controle de uma prótese de mão e de um braço robótico. & \cite{Hochberg2006} & medido \\
10 & Escrita com a ``mão mental'': cerca de $90$ caracteres por minuto, menos de $8$~bit/s & Pessoa com paralisia da mão, tentativas de escrever à mão, decodificador recorrente: $90$ caracteres por minuto com precisão de $94.1\,\%$ em tempo real, acima de $99\,\%$ após autocorreção. A estimativa em bits é aritmética do capítulo. & \cite{Willett2021} & medido \\
11 & Tentativas de falar são convertidas em texto a cerca de $60$ palavras por minuto & Pessoa que perdeu a fala inteligível, matrizes no córtex: $62$ palavras por minuto; $9.1\,\%$ de erros de palavra com vocabulário de $50$ palavras, $23.8\,\%$ com vocabulário de $125\,000$. & \cite{Willett2023} & medido \\
12 & A interface extrai uma pequena fração da capacidade: um neurônio dá dezenas de bits por segundo, cem dão milhares (proposição~\ref{prop:probe}) & O limite superior~\eqref{eq:spikeentropy} é teoria; a comparação com as linhas~8 e 10 é conclusão do capítulo. Que o gargalo sejam as poucas dimensões e o desconhecimento do código, e não o fio, não foi verificado por experimento direto. & \cite{Rieke1997} & teoria \\
13 & O cérebro e o decodificador aprendem um com o outro & Macacos controlam um cursor; o decodificador se ajusta em malha fechada, e os neurônios reaprendem ao mesmo tempo: a precisão aumenta, a habilidade se mantém e sofre menos com os movimentos do próprio braço. O conselho de separar as velocidades de aprendizado é uma regra de engenharia; o capítulo não traz um experimento específico. & \cite{Orsborn2014coad} & medido \\
14 & A corrente num eletrodo excita os neurônios e fios em volta sem distinguir tipo & Córtex de camundongo e gato, registro de cálcio por dois fótons: mesmo uma corrente fraca excita neurônios esparsos, espalhados a distâncias de até milímetros, por excitação direta de axônios num volume de dezenas de micrômetros de diâmetro. Ou seja, a excitação não é seletiva, mas também não é contínua. & \cite{Histed2009stim} & medido \\
15 & A estimulação do córtex visual acende pontos; até $16$ pontos ao mesmo tempo permitem reconhecer algumas letras e bordas de objetos & Pessoa cega, $96$ eletrodos no córtex visual por $6$ meses: limiar de um eletrodo de $66.8\pm36.5$~\textmu A; a estimulação simultânea de grupos (até $16$ eletrodos, o limite do estimulador) reduz o limiar e produz padrões distinguíveis; algumas letras e bordas de objetos são reconhecidas. & \cite{Fernandez2021} & medido \\
16 & A segunda linha da Neuralink é um dispositivo que envia sinais ao córtex visual & Só comunicados da empresa sobre o desenvolvimento; não foram encontrados dados revisados por pares sobre seu funcionamento. & relatórios da Neuralink & hipótese \\
17 & As concentrações dos neurotransmissores de difusão podem ser medidas por um sensor químico no tecido & Fatias de cérebro de rato, voltametria cíclica rápida com fibra de carbono: medidas a liberação e a recaptação de serotonina; a substância sai para além da sinapse. & \cite{Bunin1998} & medido \\
\bottomrule
\end{longtable}}

\section{Capítulo~\ref{sec:eetypes}. Neurônios como componentes}

Os modos de funcionamento de células individuais foram medidos diretamente, com corrente por eletrodo e registro de tensão; a matemática do integrador e da divisão em classes é teoria clássica, e o uso da reconfiguração de modos pelo cérebro para calcular continua sendo uma hipótese.

{\footnotesize
\setlength{\tabcolsep}{3pt}\renewcommand{\arraystretch}{1.15}
\begin{longtable}{>{\raggedright}p{0.5cm}>{\raggedright}p{4.0cm}>{\raggedright}p{5.6cm}>{\raggedright}p{2.3cm}>{\raggedright\arraybackslash}p{1.7cm}}
\toprule
N.º & Proposição & Experimento e o que foi medido & Fonte & Status \\
\midrule\endfirsthead
\toprule
N.º & Proposição & Experimento e o que foi medido & Fonte & Status \\
\midrule\endhead
1 & Ressonador: uma corrente lenta que se opõe à mudança de tensão torna o neurônio um filtro passa-faixa com máximo entre alguns hertz e dezenas de hertz (seção~\ref{sec:resonator}) & Em fatias, injetou-se em neurônios \nt{GLUT} corrente senoidal com frequência crescente e mediu-se a impedância $|Z(f)|$: máximo em $2$--$5$~Hz a $33^\circ$C (cerca de $7$~Hz a $38^\circ$C); a ressonância é criada por correntes lentas de potássio e catiônica e amplificada pela corrente persistente de sódio. Uma revisão mostra o mesmo recurso em muitas classes de células & \cite{HuVervaekeStorm2002,HutcheonYarom2000} & medido \\
2 & A corrente lenta se comporta como indutância e, com a capacitância da membrana, forma um circuito ressonante & A resposta sublimiar do axônio à corrente foi medida e descrita por equações de condutância linearizadas; o circuito equivalente contém uma indutância ``fenomenológica'' cujo valor depende da tensão & \cite{Mauro1970} & teoria \\
3 & Constante de tempo do grupo com realimentação $\tau_{\text{ef}}=\tau/(1-w)$~\eqref{eq:perfectint}; com $\tau=0.1$~s e $w=0.995$, $20$~s & Deduzida no capítulo a partir da equação linear do grupo; proposta como mecanismo de manutenção da posição do olho num trabalho teórico & dedução no capítulo; \cite{Seung1996} & teoria \\
4 & O integrador da posição do olho mantém a entrada pela rede, e não por uma propriedade de uma célula isolada & Registro intracelular em peixe, no núcleo que mantém a posição do olho: após um giro rápido, a taxa do neurônio muda em degrau e se mantém; uma corrente breve numa célula causa só um deslocamento passageiro, e os degraus de tensão persistem abaixo do limiar: são mantidos pelas entradas sinápticas & \cite{Aksay2001} & medido \\
5 & O integrador sem vazamento precisa de ajuste contínuo por aprendizado; desajustado, ele esquece ou satura & Peixes treinados com realimentação visual proporcional a $\pm$ a posição do olho: a fixação ficava instável ou com vazamento, e a constante de tempo dos neurônios caía mais de uma ordem de grandeza, para $<1$~s; a realimentação visual normal restabelecia aos poucos a estabilidade & \cite{Major2004} & medido \\
6 & Neurônio biestável: uma entrada curta liga um platô duradouro, a inibição o desliga; a propriedade é ligada pela serotonina (seção~\ref{sec:bistable}) & Registro intracelular de neurônios \nt{ACET} que controlam músculos, em gato: uma excitação sináptica curta leva o neurônio a funcionamento prolongado, uma inibição curta o reinicia; no gato espinhal, o estado aparece após a administração de um precursor da serotonina & \cite{Hounsgaard1988} & medido \\
7 & Duas classes: integradores (a taxa cresce a partir do zero) e ressonadores (no limiar, frequência finita de imediato) & As classes decorrem da teoria de bifurcações. Experimento: em fatias de córtex, neurônios \nt{GLUT} começam a disparar com taxa tão baixa quanto se queira (tipo~1), e os neurônios \nt{GABA} rápidos, com salto para uma frequência finita (tipo~2) & \cite{Izhikevich2007,Tateno2004} & medido \\
8 & Um neurônio é integrador ou detector de coincidência: a inibição encurta a janela de soma & Em fatias de hipocampo, a janela em que as entradas excitatórias se somam até o limiar é menor que $2$~ms; ela é encurtada pela inibição direta que chega logo após a excitação; nos dendritos, onde a inibição é mais fraca, a janela é mais larga & \cite{Pouille2001} & medido \\
9 & Linha de atraso feita de axônios (tabela~\ref{tab:eetypes}) & Na coruja-das-torres foram medidos os atrasos nos axônios que vão aos detectores de coincidência: comprimentos diferentes de axônio dão atrasos diferentes, e os detectores estão sintonizados com a diferença dos tempos de chegada do som & \cite{Carr1990} & medido \\
10 & Marca-passo na vigília e célula silenciosa no sono & Os neurônios \nt{SERO} funcionam de modo quase regular de dia, desaceleram no sono de ondas lentas e quase se calam no sono REM (mais detalhes no capítulo~\ref{sec:sleep}) & \cite{JacobsAzmitia1992,McGintyHarper1976} & medido \\
11 & O componente é definido pelos canais iônicos e pelos canais de difusão que os ajustam (proposição~\ref{prop:eetypes}) & Mostrado em redes pequenas: substâncias moduladoras alteram as propriedades intrínsecas dos neurônios e a força das sinapses, de modo que o mesmo circuito produz ritmos diferentes & \cite{Marder2012} & medido \\
12 & O cérebro usa a reconfiguração de modos para calcular (proposição~\ref{prop:eetypes}) & Não há experimento direto em que a reconfiguração do modo de uma célula seja necessária para resolver uma tarefa; há só experimentos em que os moduladores mudam a saída do circuito & \cite{Marder2012} & hipótese \\
\bottomrule
\end{longtable}}

\section{Capítulo~\ref{sec:dendrites}. Número de dendritos}

A estrutura das classes e o número de entradas foram medidos por anatomia e por registros elétricos; a estimativa~\eqref{eq:dendriteload} é física simples do capacitor, e a explicação computacional do número de entradas se apoia em teoria e modelos.

{\footnotesize
\setlength{\tabcolsep}{3pt}\renewcommand{\arraystretch}{1.15}
\begin{longtable}{>{\raggedright}p{0.5cm}>{\raggedright}p{4.0cm}>{\raggedright}p{5.6cm}>{\raggedright}p{2.3cm}>{\raggedright\arraybackslash}p{1.7cm}}
\toprule
N.º & Proposição & Experimento e o que foi medido & Fonte & Status \\
\midrule\endfirsthead
\toprule
N.º & Proposição & Experimento e o que foi medido & Fonte & Status \\
\midrule\endhead
1 & Capacitância específica da membrana $c_m\approx1$~\textmu F/cm$^2$ & Retirava-se membrana do corpo do neurônio com uma pipeta, aplicava-se um degrau de tensão e, pelo decaimento da corrente capacitiva, achava-se a capacitância total, depois dividida pela área: $0.9$~\textmu F/cm$^2$ em três classes de neurônios & \cite{Gentet2000} & medido \\
2 & Tempo de carga $t\approx c_mA\Delta V/I$~\eqref{eq:dendriteload}: o neurônio grande precisa de dezenas de entradas simultâneas, ao pequeno basta uma sinapse grande & Estimativa pela lei de carga do capacitor; os números ($20$ e $300$~pF, $0.1$ e alguns nA) são ordens de grandeza & dedução no capítulo & teoria \\
3 & Uma sinapse enorme dá corrente de alguns nanoampères e leva o neurônio pequeno ao limiar de imediato & Registro simultâneo do terminal e do destinatário numa fatia: corrente de uma sinapse de $2$--$13$~nA, cada disparo de entrada provoca resposta supralimiar, atraso de cerca de $0.4$~ms a $36^\circ$C; o destinatário emite \nt{GLYC} & \cite{Borst1995,BorstSoria2012} & medido \\
4 & Zero dendritos: nos neurônios sensoriais do tato e da dor, o disparo vai do sensor à medula espinhal sem passar pelo corpo celular & A estrutura (axônio que se divide em dois junto ao corpo) foi estabelecida anatomicamente. Que a condução não exige excitabilidade do corpo foi mostrado num modelo validado desse neurônio: sem corpo excitável, o disparo atravessa a bifurcação com a mesma confiabilidade & \cite{AmirDevor2003} & teoria \\
5 & Um dendrito: o neurônio olfatório carrega um tipo de receptor e transmite uma linha de entrada & Revisão de experimentos com genes de receptores: cada neurônio olfatório expressa um gene de receptor de uma grande família & \cite{Mombaerts2004} & medido \\
6 & Dois dendritos: nos detectores de direção do som, as entradas do ouvido esquerdo e do direito chegam a dendritos diferentes & Anatomia e registros em mamíferos, reunidos numa revisão: as entradas dos dois ouvidos são separadas em dois dendritos opostos & \cite{Grothe2010} & medido \\
7 & Separar as entradas em dois dendritos torna o ``E'' mais estrito & Mostrado em modelo: a saturação~\eqref{eq:shunt} num só dendrito impede que as entradas de um lado cheguem ao limiar, e a detecção de coincidências melhora. Não há experimento direto em que se tenha mudado a distribuição das entradas & \cite{AgmonSnir1998} & teoria \\
8 & Cerca de $7\cdot10^{10}$ pequenos neurônios \nt{GLUT} do circuito de aprendizado motor, com $\sim4$ entradas cada & Contagem de células por fracionamento isotrópico: nessa parte do cérebro humano há cerca de $69$ bilhões de neurônios. Registro de uma dessas células em rato vivo: ao toque chegam rajadas de correntes discretas de poucas entradas, e a célula dispara quando elas se somam & \cite{Azevedo2009,Chadderton2004} & medido \\
9 & Um elemento de limiar com $n$ entradas separa no máximo $P_{\max}\approx2n$ padrões aleatórios~\eqref{eq:cover}; para o neurônio de saída do circuito de aprendizado motor, o limite é $\sim2\cdot10^5$, e a estimativa realista, $\sim4\cdot10^4$ associações & O limite é um resultado clássico da geometria combinatória. A estimativa realista vem da teoria de capacidade com pesos só positivos e margem para o ruído, aplicada ao número medido de entradas desse neurônio & \cite{Cover1965,Brunel2004} & teoria \\
10 & Os misturadores de poucas entradas servem para expandir a entrada; ``quatro'' está perto do ótimo & Teoria: a dimensão da representação dada pela camada de misturadores é máxima aproximadamente com o número de entradas observado na anatomia; a hipótese original da expansão está nos trabalhos clássicos & \cite{Marr1969,Albus1971,LitwinKumar2017} & hipótese \\
11 & Árvores grandes: $10^4$--$10^5$ entradas & Contagem de sinapses em imagens de microscopia eletrônica: $\approx1.7\cdot10^5$ sinapses num neurônio \nt{GABA} de saída do circuito de aprendizado motor em rato; cerca de $30\,000$ entradas excitatórias e $1\,700$ inibitórias num neurônio \nt{GLUT} do hipocampo & \cite{NapperHarvey1988,Megias2001} & medido \\
12 & Os ramos de uma árvore grande são subcircuitos separados com limiares próprios; o neurônio é uma pequena rede de várias camadas & Estimulação pontual de dois locais em dendritos finos: entradas no mesmo ramo se somam de forma sigmoide, em ramos diferentes, linearmente. Nos dendritos de neurônios do córtex humano foram encontrados disparos graduados de cálcio que permitem a uma só célula resolver uma tarefa linearmente não separável. Modelo: para reproduzir um neurônio é preciso uma rede profunda & \cite{Polsky2004,Gidon2020,Beniaguev2021} & medido \\
13 & Dendritos-saída: cada ramo do neurônio \nt{GABA} da retina é um detector de direção separado & Registro de cálcio por dois fótons em ramos isolados: o sinal no ramo é seletivo ao movimento do corpo celular para a ponta, e a tensão do corpo não & \cite{Euler2002} & medido \\
14 & O número de entradas segue a tarefa (proposição~\ref{prop:dendrites}) & A estrutura das classes foi medida (linhas 3--13); que o número de entradas foi escolhido pela velocidade ou pela classificação só foi mostrado por teoria e modelos, para algumas classes & \cite{LitwinKumar2017,AgmonSnir1998} & hipótese \\
\bottomrule
\end{longtable}}

\section{Capítulo~\ref{sec:sleep}. Sono}

Os modos do sono, os replays e a redução das sinapses foram medidos por registros de neurônios, microdiálise e microscopia eletrônica; o horário do sono e a gangorra lenta são descritos pelos modelos do livro, e a finalidade do sono é, em grande parte, hipótese.

{\footnotesize
\setlength{\tabcolsep}{3pt}\renewcommand{\arraystretch}{1.15}
\begin{longtable}{>{\raggedright}p{0.5cm}>{\raggedright}p{4.0cm}>{\raggedright}p{5.6cm}>{\raggedright}p{2.3cm}>{\raggedright\arraybackslash}p{1.7cm}}
\toprule
N.º & Proposição & Experimento e o que foi medido & Fonte & Status \\
\midrule\endfirsthead
\toprule
N.º & Proposição & Experimento e o que foi medido & Fonte & Status \\
\midrule\endhead
1 & No sono os neurônios do córtex não se calam; o que muda é o modo de funcionamento & Registros longos de neurônios do córtex em ratos: no sono de ondas lentas a taxa continua não nula, e períodos de atividade se alternam com períodos de silêncio; após vigília longa a taxa é maior em todos os estados, após o sono, menor & \cite{Vyazovskiy2009} & medido \\
2 & O \nt{ACET} é alto de dia e no sono REM, baixo no sono de ondas lentas (tabela~\ref{tab:sleepmodes}) & Microdiálise no córtex de gatos em livre movimento: na vigília, a liberação de acetilcolina é $100$--$175\%$ maior que no sono de ondas lentas, e no sono REM fica mais ou menos no nível da vigília tranquila & \cite{Marrosu1995} & medido \\
3 & \nt{NORA} e \nt{SERO} se aquietam no sono de ondas lentas e quase se calam no REM & Registro de neurônios \nt{NORA} isolados em ratos e de neurônios \nt{SERO} em gatos em cada estágio do sono: a taxa é máxima na vigília, menor no sono de ondas lentas e quase nula no REM & \cite{AstonJonesBloom1981,McGintyHarper1976} & medido \\
4 & No sono REM os músculos são desligados por inibição via \nt{GABA} e \nt{GLYC} & Em ratos, mediu-se a atividade dos neurônios \nt{ACET} do músculo masseter e aplicaram-se bloqueadores diretamente sobre eles: a paralisia só cessa se forem bloqueados tanto os receptores \nt{GABA} dos dois tipos quanto os receptores \nt{GLYC} & \cite{BrooksPeever2012} & medido \\
5 & A pressão do sono $S$ é um capacitor com dois limiares~\eqref{eq:sleeppressure}, $\tau_r=18.2$~h, $\tau_d=4.2$~h & Modelo de dois processos; as constantes de tempo foram obtidas de como a potência das ondas lentas do EEG cresce com a vigília e cai no sono, e a forma dos limiares, do horário do despertar espontâneo & \cite{Borbely1982,Daan1984} & teoria \\
6 & A máquina chega sozinha a $16.1$~h de vigília e $8.0$~h de sono (fig.~\ref{fig:twoprocess}) & Cálculo por~\eqref{eq:sleeppressure} com limiares oscilantes, script \texttt{models/sleep\_models.py}: $16.1$~h de vigília e $7.9$--$8.0$~h de sono & --- & modelo \\
7 & Após $40$~h sem dormir, $S=0.90$, mas o sono dura só $8.8$~h: a dívida é paga com profundidade, não com duração & Modelo: \texttt{models/sleep\_models.py} dá $S=0.90$ e $8.8$~h. Experimento: após $40.5$~h sem dormir, em oito pessoas, na primeira noite de recuperação, a potência das ondas lentas do EEG foi significativamente maior que o normal & \cite{Borbely1981} & modelo \\
8 & Fisicamente, $S$ é a adenosina & Microdiálise em gatos: a adenosina extracelular numa das regiões que controlam a vigília cresce durante a vigília, cresce ainda mais com a privação de sono e cai durante o sono. Que $S$ seja só adenosina não foi mostrado & \cite{PorkkaHeiskanen1997,Dunwiddie2001} & hipótese \\
9 & No sono de ondas lentas o córtex oscila: atividade por algumas centenas de ms, depois silêncio, menos de uma vez por segundo ($<1$~Hz, sob anestesia em geral $0.3$--$0.4$~Hz) & Registro intracelular em gatos anestesiados: despolarização de $0.8$--$1.5$~s e hiperpolarização longa, ritmo $<1$~Hz, em geral $0.3$--$0.4$~Hz; a mesma alternância de atividade e silêncio aparece no sono natural (linha~1) & \cite{Steriade1993,Vyazovskiy2009} & medido \\
10 & A gangorra é um grupo com realimentação e fadiga~\eqref{eq:updown}; com $b=5$, período de $1.1$~s (valor do modelo) e atividade em cerca de $35\%$ do tempo; com $b=1$, atividade contínua (fig.~\ref{fig:updown}) & Cálculo, script \texttt{models/sleep\_models.py}: período de $1.11$~s, fração de tempo ativo $0.35$ (média sobre um número inteiro de períodos); com $b=1$, fração ativa de $1.0$. O período do modelo é mais curto que o medido (linha~9) & --- & modelo \\
11 & O \nt{ACET} desliga a gangorra e leva o córtex ao funcionamento contínuo & A estimulação de um núcleo de neurônios \nt{ACET} em gato bloqueou o ritmo lento em $79\%$ dos neurônios do córtex e os levou a funcionamento contínuo, sobretudo pelo desaparecimento das hiperpolarizações longas. A acetilcolina suprime as correntes lentas de potássio que criam a fadiga & \cite{SteriadeAmzica1993,McCormick1992} & medido \\
12 & Os surtos de ritmo de $7$--$15$~Hz são gerados pela malha \nt{GLUT}--\nt{GABA} entre o córtex e um nó de retransmissão & Em fatias do nó de retransmissão visual de furão, os surtos são criados por rajadas de resposta das células de retransmissão à inibição vinda de um núcleo \nt{GABA} vizinho, que elas mesmas excitam & \cite{vonKrosigk1993,SteriadeMcCormickSejnowski1993} & medido \\
13 & Os replays do dia são acelerados: no hipocampo cerca de $20$ vezes, no córtex $6$--$7$ vezes & No sono de ondas lentas, as sequências repetidas de neurônios do hipocampo aparecem comprimidas no tempo. Depois de correr numa pista, as sequências de neurônios de lugar se repetiram na mesma ordem em surtos curtos de $\sim100$~ms, comprimidas cerca de $20$ vezes; no córtex, compressão de $6$--$7$ vezes & \cite{Nadasdy1999,LeeWilson2002,Euston2007} & medido \\
14 & Reforçar a coordenação dos replays com o córtex melhora a memória (relação causal) & Em ratos, após o treino, uma estimulação sincronizada com os replays reforçou sua ligação com as ondas lentas e os surtos de ritmo no córtex: no dia seguinte a memória era alta, e nos controles, no nível do acaso & \cite{Maingret2016} & medido \\
15 & A finalidade dos replays é o aprendizado intercalado da memória lenta & Teoria dos dois sistemas de aprendizado e cálculos em redes artificiais: o aprendizado sequencial estraga o antigo, o intercalado não. Não há experimento direto no cérebro mostrando que os replays servem justamente a isso & \cite{McClelland1995} & hipótese \\
16 & Após o sono, as sinapses diminuem proporcionalmente ao tamanho, e as grandes quase não mudam & Microscopia eletrônica tridimensional de $6920$ sinapses do córtex de camundongo: a área de contato após o sono é $\sim18\%$ menor que após a vigília, a redução é proporcional ao tamanho, e as sinapses grandes e estáveis não mudam & \cite{deVivo2017} & medido \\
17 & A tarefa principal do sono é a renormalização dos pesos & Hipótese da homeostase sináptica; as linhas 1 e 16 concordam com ela, mas não provam que seja a função principal & \cite{TononiCirelli2014} & hipótese \\
18 & O sono REM é um teste em bancada do modelo do mundo & O modo (tabela~\ref{tab:sleepmodes}) foi medido; que a rede execute o modelo do mundo é uma suposição teórica, sem experimento & \cite{Hobson2009} & hipótese \\
19 & Lavagem do cérebro no sono: questão em aberto & Um experimento: no sono, o espaço extracelular é $60\%$ mais largo e a limpeza, mais rápida. Outro, com outro método de medida: no sono e sob anestesia, a limpeza é bem mais lenta. Os resultados se contradizem & \cite{Xie2013,Miao2024} & hipótese \\
20 & Sono local: trechos do córtex de um animal acordado caem por instantes no silêncio & Em ratos, após vigília longa, os neurônios de um trecho isolado do córtex se calam brevemente, como no sono de ondas lentas, com uma onda lenta no EEG local; nesses momentos o rato erra mais na tarefa & \cite{Vyazovskiy2011} & medido \\
\bottomrule
\end{longtable}}

\section{Capítulo~\ref{sec:livingmachine}. A máquina de neurônios vivos}

O comportamento das culturas em matrizes de eletrodos foi medido em experimentos diretos com registro e estimulação; o mecanismo das salvas se apoia no modelo de esgotamento das sinapses e em experimentos com microculturas, e as afirmações sobre a dopamina na placa, em experimentos isolados, parte dos quais os próprios autores consideram apenas uma demonstração.

{\footnotesize
\setlength{\tabcolsep}{3pt}\renewcommand{\arraystretch}{1.15}
\begin{longtable}{>{\raggedright}p{0.5cm}>{\raggedright}p{4.0cm}>{\raggedright}p{5.6cm}>{\raggedright}p{2.3cm}>{\raggedright\arraybackslash}p{1.7cm}}
\toprule
N.º & Proposição & Experimento e o que foi medido & Fonte & Status \\
\midrule\endfirsthead
\toprule
N.º & Proposição & Experimento e o que foi medido & Fonte & Status \\
\midrule\endhead
1 & Cultura no chip: sobretudo \nt{GLUT}, com uma fração menor de \nt{GABA} que depende da preparação; em culturas de hipocampo, cerca de $6\%$ & Uma rede de milhares de neurônios sobre uma matriz de eletrodos é um experimento padrão (linha~3). A fração de neurônios \nt{GABA} em culturas de hipocampo foi contada por marcação para GABA: cerca de $6\%$ dos neurônios. Para culturas de córtex não damos um número verificado & \cite{Benson1994} & medido \\
2 & Organoides: tecido nervoso tridimensional com cerca de um milímetro, a partir de células-tronco humanas & A partir de células-tronco foram cultivados organoides tridimensionais com várias regiões cerebrais, entre elas um córtex com camadas de precursores e neurônios maduros & \cite{Lancaster2013} & medido \\
3 & Sem entrada, a cultura dispara em salvas com intervalos de um segundo a minutos, conforme a cultura & $58$ culturas de córtex com densidade de $3\,000$ a $50\,000$ neurônios foram registradas por cinco semanas ($963$ registros de meia hora): depois de duas semanas, em quase todas as culturas predominam salvas da rede inteira; os intervalos entre salvas vão de $1$ a $300$~s e dependem fortemente da preparação & \cite{Wagenaar2006} & medido \\
4 & A salva termina por esgotamento do neurotransmissor, intervalo $T_{\text{salva}}\approx\tau_{\text{rec}}\ln\frac{1}{1-x^*}$~\eqref{eq:burstinterval}; $2$--$4$~s é a estimativa do modelo com $\tau_{\text{rec}}=0.8$~s, a recuperação medida é mais lenta & A fórmula é deduzida no capítulo. Modelo: uma rede com sinapses que se esgotam gera sozinha salvas da rede inteira. Experimento em microculturas de $4$--$30$ neurônios: a duração da salva depende do tempo desde a anterior, e o esgotamento das vesículas de neurotransmissor elimina as salvas; conclusão dos autores: a salva é limitada pelo estoque de neurotransmissor. Lá a recuperação do estoque leva dezenas de segundos, o que, com a mesma fórmula, dá intervalos de dezenas de segundos e minutos & dedução no capítulo; \cite{Tsodyks2000,CohenSegal2011,Tsodyks1997} & teoria \\
5 & ``O professor que se cala'': a estimulação é desligada quando aparece a resposta desejada, e a resposta se fixa & A cultura foi estimulada a $0.3$--$1$~Hz até aparecer a resposta escolhida $50\pm10$~ms depois do estímulo; então a estimulação era desligada por $5$~min; depois de alguns ciclos a resposta desejada surgia de imediato; foram traçadas curvas de aprendizado & \cite{Shahaf2001} & medido \\
6 & A cultura joga tênis; aumento significativo das sequências na sessão só com realimentação completa & Em detalhe no capítulo~\ref{sec:dishbrain} e em seu suplemento: com silêncio no lugar do estímulo, a cultura também joga melhor no fim da sessão do que sem realimentação, mas com efeito menor; o aprendizado de uma sessão para outra é instável & \cite{Kagan2022} & medido \\
7 & Que a cultura aprende a prever sua entrada é hipótese; as palavras fortes sobre ``inteligência'' foram contestadas & O princípio da energia livre é uma proposta teórica; não há experimento direto que o distinga de explicações mais simples. Trinta pesquisadores, num artigo de resposta, apontaram que a mudança de comportamento na malha não mostra nem inteligência nem sensações & \cite{Friston2010,Balci2023} & hipótese \\
8 & A cultura conduz um corpo virtual até o alvo com um padrão de estímulos escolhido & O programa escolhia sozinho os padrões de estimulação de treino pelo resultado corrente; em dezenas de minutos a rede atingia os estados pedidos, e a plasticidade durava $80$~min após $2$~h de treino. Os mesmos estímulos, aplicados sem relação com o comportamento, não tinham efeito & \cite{Bakkum2008} & medido \\
9 & Reservatório: só se treina a leitura linear $y=W_{\text{saída}}\mathbf r$~\eqref{eq:reservoir} & Teoria da computação de reservatório: qualquer sistema não linear com memória que se apaga, mais uma saída linear treinada & \cite{Maass2002,JaegerHaas2004} & teoria \\
10 & O organoide como reservatório distingue falantes com acurácia de cerca de $78\%$ (segundo os autores); a leitura aprende pelo erro, e as conexões do organoide mudam sem professor & Sons de fala de vários falantes foram aplicados ao organoide como padrões de corrente na matriz de eletrodos; a leitura treinada distinguia o falante com acurácia de cerca de $78\%$, segundo os autores. Eles também relatam que durante o treino a conectividade funcional do organoide mudou, e chamam isso de aprendizado não supervisionado no próprio organoide & \cite{Cai2023} & medido \\
11 & Um milhão de neurônios gasta cerca de $10^{-4}$~W em disparos~\eqref{eq:dishpower} & Estimativa: o custo de um disparo, $10^{-10}$~J, do orçamento energético do córtex~\eqref{eq:energy}, multiplicado pelo número de disparos; não há medição direta da potência da cultura & dedução no capítulo; \cite{Attwell2001} & teoria \\
12 & Dopamina por clarão de luz: $1$~mM de dopamina enjaulada, $240$ repetições, o número de disparos evocados cresceu & Numa plataforma remota de organoides, dopamina numa gaiola fotossensível ($1$~mM) era liberada com ultravioleta ($365$~nm, $0.8$~s) quando o estímulo evocava mais disparos; em $240$ passos (cerca de $5$~h) o número de disparos em geral cresceu. Os autores escrevem que um único experimento não permite estabelecer a relação com a dopamina; não há controles & \cite{Jordan2024} & hipótese \\
13 & A dopamina de células vivas muda o peso de um transistor orgânico de forma duradoura e reversível & Células que liberam dopamina foram cultivadas sobre um elemento neuromórfico orgânico; a oxidação da dopamina junto ao eletrodo muda a condutância do canal; foram mostrados a mudança duradoura do peso e seu retorno & \cite{Keene2020} & medido \\
14 & Existem organoides com neurônios \nt{DOPA} funcionais; sua dopamina não foi usada como professor & Em organoides de células-tronco humanas foram encontrados neurônios \nt{DOPA} maduros e eletricamente ativos, e produção de dopamina. Não encontramos na literatura experimentos em que essa dopamina ensine outra rede & \cite{Jo2016} & medido \\
15 & O organoide equilibra a haste: os estímulos de ensino do programa são melhores que os aleatórios, sem consolidação, por meio das sinapses \nt{GLUT} & Organoides de córtex de camundongo em malha fechada com a tarefa do ``pêndulo no carrinho'': na maioria dos organoides, os sinais escolhidos por aprendizado por reforço deram resultado melhor que os aleatórios ou nenhum; após $45$~min de repouso a melhora não se mantinha; o bloqueio dos receptores AMPA e NMDA a anulava & \cite{Robbins2026} & medido \\
16 & Sem registradores de controle, a rede na bancada não consegue decidir sozinha o que conta como sucesso (proposição~\ref{prop:livingmachine}) & Em todos os experimentos das linhas 5--15, o sinal de ensino é definido pelo experimentador ou por um programa; não há experimento em que a cultura tenha criado sozinha um critério de sucesso; é uma conclusão pela ausência, não um experimento & --- & hipótese \\
\bottomrule
\end{longtable}}

\section{Capítulo~\ref{sec:dishbrain}. DishBrain}

A montagem da bancada e os resultados do jogo vêm de um único trabalho experimental e de sua continuação; as fórmulas do ruído e da deriva para boas configurações são deduzidas no capítulo e verificadas com o modelo do livro, e o mecanismo de aprendizado na placa não está estabelecido.

{\footnotesize
\setlength{\tabcolsep}{3pt}\renewcommand{\arraystretch}{1.15}
\begin{longtable}{>{\raggedright}p{0.5cm}>{\raggedright}p{4.0cm}>{\raggedright}p{5.6cm}>{\raggedright}p{2.3cm}>{\raggedright\arraybackslash}p{1.7cm}}
\toprule
N.º & Proposição & Experimento e o que foi medido & Fonte & Status \\
\midrule\endfirsthead
\toprule
N.º & Proposição & Experimento e o que foi medido & Fonte & Status \\
\midrule\endhead
1 & Bancada: cerca de $800\,000$ células de córtex de camundongo ou cerca de $10^6$ células humanas por chip; $26\,000$ eletrodos em $8$~mm$^2$, registro de até $1024$, estimulação por $8$ & Métodos do trabalho: chip MaxOne, $26\,000$ eletrodos de platina em $8$~mm$^2$, registro de $1024$ canais a $20\,000$ amostras por segundo; tecnicamente é possível estimular até $32$ eletrodos, de modo independente conseguiram $8$; as culturas de camundongo foram usadas a partir de cerca de $10$~dias & \cite{Kagan2022} & medido \\
2 & Malha com ciclo de $10$~ms: os disparos das regiões motoras movem a raquete (fig.~\ref{fig:dishloop}) & Os disparos são contados em janelas de $10$~ms; o atraso entre o disparo e o estímulo de resposta é de cerca de $5$~ms, definido pelo buffer do equipamento & \cite{Kagan2022} & medido \\
3 & Entrada: código de lugar (qual dos $8$ eletrodos) e de taxa, $4$--$40$~disparos/s & No experimento principal, a frequência de estimulação cresce linearmente de $4$~Hz, com a bola junto à parede do fundo, até $40$~Hz junto à raquete; a amplitude dos pulsos da bola é de $75$~mV; os pulsos são retangulares bifásicos & \cite{Kagan2022} & medido \\
4 & Saída $\Delta p=c\,(n_\uparrow-n_\downarrow)$~\eqref{eq:dishreadout}, ruído da contagem por janela $1/\sqrt{RT}$ & A regra da diferença entre duas regiões é descrita no trabalho (com pesos das regiões pela atividade), a fórmula exata não é dada. A estimativa do ruído é consequência da contagem de Poisson, deduzida no capítulo & \cite{Kagan2022}; dedução no capítulo & teoria \\
5 & Professor: após o acerto, estímulo previsível de $75$~mV, $100$~disparos/s, $0.1$~s; após o erro, imprevisível, $150$~mV, $5$~disparos/s, $4$~s em lugares e momentos aleatórios, depois $4$~s de pausa & Métodos: após o acerto, $75$~mV, $100$~Hz, $100$~ms em todos os $8$ eletrodos de uma vez; após o erro, $150$~mV, $5$~Hz em eletrodos aleatórios e momentos aleatórios durante $4$~s, depois $4$~s sem estimulação. Não há dopamina na cultura & \cite{Kagan2022} & medido \\
6 & A rede age de modo a tornar sua entrada previsível & O princípio da energia livre é uma proposta teórica da qual os autores derivaram o protocolo; o experimento é compatível com ele, mas não o distingue de mecanismos mais simples (linhas~7--8) & \cite{Friston2010,Kagan2022} & hipótese \\
7 & Se o fracasso sacode a rede e o sucesso não, o tempo numa configuração é $\rho\propto1/P_{\text{erro}}$~\eqref{eq:dishdensity} (proposição~\ref{prop:dishscramble}) & Decorre do balanço de fluxos numa cadeia de Markov com empurrões simétricos, deduzido no capítulo. Não há teste direto numa cultura & dedução no capítulo & teoria \\
8 & O modelo ``embaralhar no erro'' aprende: fração de acertos $0.21\to0.38$; com empurrão após cada rali $\approx0.12$, sem empurrões $0.19$ (fig.~\ref{fig:dishmodel}) & Cálculo, script \texttt{models/dishbrain\_model.py}, $40$ culturas-modelo com $2000$ ralis cada: $0.21\to0.38$; $0.13\to0.11$; $0.19\to0.19$ & --- & modelo \\
9 & As sequências de bolas rebatidas se alongam só nas culturas vivas na malha completa; os controles não aprendem & $399$ sessões de $20$~min: $9$ culturas de camundongo ($101$ sessões), $11$ humanas ($138$), $6$ chips com meio ($80$), $20$ culturas em repouso ($42$), $3$ simulações ($38$). A duração média do rali nos últimos $15$~min é maior que nos primeiros $5$ só nas culturas vivas; nas células que não são neurônios (HEK293T) e nos chips com meio não há efeito & \cite{Kagan2022} & medido \\
10 & Aumento significativo na sessão só com realimentação completa; com silêncio no lugar do estímulo, o jogo no fim da sessão é melhor que sem realimentação, mas o efeito é menor & $486$ sessões em $3$ dias: ``estímulo'' ($n=27$), ``silêncio'' ($17$), ``sem realimentação'' ($15$). O aumento ao longo do tempo só é significativo na condição ``estímulo''; no fim da sessão, ``silêncio'' superava ``sem realimentação'' com efeito menor & \cite{Kagan2022} & medido \\
11 & O efeito é pequeno; se a rede aprende de forma duradoura é questão em aberto & Dentro da sessão a melhora é sólida, mas, segundo os próprios autores, um aprendizado estável entre sessões ao longo de vários dias não foi observado & \cite{Kagan2022} & medido \\
12 & Durante o jogo a rede se aproxima do regime crítico, e quanto mais perto, melhor o jogo & Análise de avalanches de atividade nas mesmas culturas: a entrada estruturada aproxima a rede do regime crítico, e a qualidade do jogo se correlaciona com essa proximidade; mas a criticidade sozinha, sem informação sobre as consequências das ações, não basta para aprender & \cite{Habibollahi2023} & medido \\
13 & A ``senciência'' da cultura não foi demonstrada & Artigo de resposta de trinta pesquisadores: a mudança de comportamento numa malha fechada não prova nem inteligência nem sensações; não há experimento que o mostre & \cite{Balci2023} & hipótese \\
14 & Quarto controle proposto: estímulo previsível após o erro, com a mesma duração e energia (seção~\ref{sec:dishspec}) & Esse experimento não foi feito; é uma proposta do livro & --- & hipótese \\
\bottomrule
\end{longtable}}

